%%%%%%%%%%%%%%%%%%%%%%%%%%%%%%%%%%%%%%%%%%%%%%%%%%%%%%%%%%%%%%%%%%%%%%%%%%%%%%%
%%  TKS_v7.4_Math.tex
%%  Extended Transfinite Fractals: Large Ordinals, Multi-Indexing, Spectra,
%%  and Functorial Constructions
%%
%%  Part of TKS v7.4 - December 2025
%%  ADDITIVE to v7.3 canon
%%%%%%%%%%%%%%%%%%%%%%%%%%%%%%%%%%%%%%%%%%%%%%%%%%%%%%%%%%%%%%%%%%%%%%%%%%%%%%%

%% ============================================================================
%% CHAPTER 1: LARGE ORDINAL CLASSES FOR NOETIC FRACTALS
%% ============================================================================

\chapter{Large Ordinal Classes for Noetic Fractals}
\label{ch:large-ordinals}

\begin{tcolorbox}[colback=blue!5!white,colframe=blue!75!black,title=\vnewfour]
This chapter extends transfinite fractal theory to incorporate large cardinals,
particularly inaccessible cardinals, which provide natural ``ceiling'' points
in the ordinal hierarchy where fractal behavior exhibits qualitative transitions.
\end{tcolorbox}

\section{Inaccessible Cardinals in Fractal Indexing}

\begin{definition}[Inaccessible Cardinal]
\label{def:inaccessible}
A cardinal $\kappa$ is \emph{inaccessible} if:
\begin{enumerate}[label=(\alph*)]
    \item $\kappa > \omega$ (uncountable),
    \item $\kappa$ is a strong limit cardinal: for all $\lambda < \kappa$,
          we have $2^\lambda < \kappa$,
    \item $\kappa$ is regular: $\mathrm{cf}(\kappa) = \kappa$.
\end{enumerate}
\end{definition}

\begin{definition}[Large Ordinal Class]
\label{def:large-ord}
The \emph{large ordinal class} $\LargeOrd$ is defined as:
\[
    \LargeOrd \;:=\; \{\, \alpha \in \Ord \mid
    \exists\, \kappa \text{ inaccessible}: \alpha < \kappa \,\}
\]
Under the assumption of at least one inaccessible cardinal, $\LargeOrd$ is a
proper class. We write $\LargeOrd_\kappa := \{\alpha < \kappa\}$ for the set
of ordinals below a specific inaccessible $\kappa$.
\end{definition}

\begin{definition}[Inaccessible-Bounded Fractal]
\label{def:inacc-fractal}
An \emph{inaccessible-bounded fractal} is a family
$\{\TransFrac{\alpha}\}_{\alpha < \kappa}$ where:
\begin{enumerate}[label=(\roman*)]
    \item Each $\TransFrac{\alpha}$ is a noetic fractal in the sense of v7.3,
    \item The family satisfies the \emph{bounded refinement property}:
    \[
        \forall\, \alpha < \beta < \kappa: \quad
        \TransFrac{\alpha} \sqsubseteq_{\mathrm{ref}} \TransFrac{\beta}
    \]
    where $\sqsubseteq_{\mathrm{ref}}$ denotes fractal refinement.
\end{enumerate}
\end{definition}

\begin{theorem}[Inaccessible Closure]
\label{thm:inacc-closure}
Let $\kappa$ be inaccessible and let $\{\TransFrac{\alpha}\}_{\alpha < \kappa}$
be an inaccessible-bounded fractal family. Then the colimit
\[
    \TransFrac{\kappa} \;:=\; \varinjlim_{\alpha < \kappa} \TransFrac{\alpha}
\]
exists in the category of noetic spaces and inherits a canonical fractal
structure.
\end{theorem}

\begin{proof}
By the regularity of $\kappa$, every cofinal subset of $\kappa$ has
cardinality $\kappa$. This ensures that the directed system
$\{\TransFrac{\alpha}\}_{\alpha < \kappa}$ satisfies the conditions for
directed colimits in the category of dcpos with Scott-continuous maps.

Since $\kappa$ is a strong limit, for any $\alpha < \kappa$ the cardinality
$|\TransFrac{\alpha}|$ satisfies $|\TransFrac{\alpha}| < \kappa$. The colimit
therefore has cardinality at most $\kappa$.

The fractal structure on $\TransFrac{\kappa}$ is induced by the universal
property: the self-similarity maps $\sigma_\alpha: \TransFrac{\alpha} \to
\TransFrac{\alpha}$ extend uniquely to $\sigma_\kappa$ by commutativity:
\[
\begin{tikzcd}
    \TransFrac{\alpha} \arrow[r, "\sigma_\alpha"] \arrow[d, hook] &
    \TransFrac{\alpha} \arrow[d, hook] \\
    \TransFrac{\kappa} \arrow[r, dashed, "\sigma_\kappa"] &
    \TransFrac{\kappa}
\end{tikzcd}
\]
\end{proof}

\begin{corollary}[Inaccessible Dimension Bound]
\label{cor:inacc-dim}
For an inaccessible-bounded fractal family:
\[
    \HausdorffDim(\TransFrac{\kappa}) \;=\; \sup_{\alpha < \kappa}
    \HausdorffDim(\TransFrac{\alpha})
\]
\end{corollary}

\section{The Ordinal Hierarchy $\OrdHierarchy$}

\begin{definition}[Ordinal Hierarchy]
\label{def:ord-hierarchy}
The \emph{ordinal hierarchy} $\OrdHierarchy$ is the structure
\[
    \OrdHierarchy \;:=\; \bigl(\Ord,\, \leq,\, \{\kappa_\xi\}_{\xi \in I},\,
    \mathrm{Lim},\, \mathrm{Succ}\bigr)
\]
where:
\begin{itemize}
    \item $\leq$ is the standard ordinal ordering,
    \item $\{\kappa_\xi\}_{\xi \in I}$ is the class of inaccessible cardinals,
    \item $\mathrm{Lim} \subset \Ord$ is the class of limit ordinals,
    \item $\mathrm{Succ}: \Ord \to \Ord$ is the successor function.
\end{itemize}
\end{definition}

\begin{definition}[Hierarchy Strata]
\label{def:strata}
For each inaccessible $\kappa_\xi$, the \emph{$\xi$-stratum} is:
\[
    \mathcal{S}_\xi \;:=\;
    \begin{cases}
        \LargeOrd_{\kappa_0} & \text{if } \xi = 0 \\
        \{\alpha \mid \kappa_{\xi-1} \leq \alpha < \kappa_\xi\} &
        \text{if } \xi \text{ is a successor} \\
        \{\alpha \mid \sup_{\zeta < \xi} \kappa_\zeta \leq \alpha < \kappa_\xi\}
        & \text{if } \xi \text{ is a limit}
    \end{cases}
\]
\end{definition}

\begin{theorem}[Stratified Fractal Decomposition]
\label{thm:stratified-decomp}
Let $\Phi^{\mathrm{trans}}$ be a transfinite fractal indexed over all of $\Ord$. Then
$\Phi^{\mathrm{trans}}$ admits a canonical decomposition:
\[
    \Phi^{\mathrm{trans}} \;\cong\; \bigoplus_{\xi \in I} \Phi^{\mathrm{trans}}|_{\mathcal{S}_\xi}
\]
where $\Phi^{\mathrm{trans}}|_{\mathcal{S}_\xi}$ denotes the restriction to stratum $\xi$.
\end{theorem}

\section{Limit Ordinal Fractal Convergence}

\begin{definition}[Limit Ordinal Fractal]
\label{def:limit-fractal}
For a limit ordinal $\lambda \in \mathrm{Lim}$, the \emph{limit ordinal fractal}
$\TransFrac{\lambda}$ is defined as:
\[
    \TransFrac{\lambda} \;:=\; \varinjlim_{\alpha < \lambda} \TransFrac{\alpha}
\]
equipped with the colimit topology and the induced self-similarity structure.
\end{definition}

\begin{lemma}[Cofinal Convergence]
\label{lem:cofinal-conv}
Let $\lambda$ be a limit ordinal with $\mathrm{cf}(\lambda) = \kappa$. Then
for any cofinal sequence $(\alpha_\nu)_{\nu < \kappa}$ in $\lambda$:
\[
    \TransFrac{\lambda} \;\cong\; \varinjlim_{\nu < \kappa} \TransFrac{\alpha_\nu}
\]
\end{lemma}

\begin{theorem}[Continuity of Hausdorff Dimension at Limits]
\label{thm:hausdorff-continuity}
The Hausdorff dimension function $\alpha \mapsto \HausdorffDim(\TransFrac{\alpha})$
is continuous at limit ordinals:
\[
    \HausdorffDim(\TransFrac{\lambda}) \;=\;
    \sup_{\alpha < \lambda} \HausdorffDim(\TransFrac{\alpha})
    \;=\; \lim_{\alpha \to \lambda^-} \HausdorffDim(\TransFrac{\alpha})
\]
\end{theorem}

\begin{proof}
The monotonicity of Hausdorff dimension under fractal refinement gives:
\[
    \alpha < \beta \implies \HausdorffDim(\TransFrac{\alpha}) \leq
    \HausdorffDim(\TransFrac{\beta})
\]
The supremum therefore exists. For the colimit $\TransFrac{\lambda}$, any
covering restricts to some $\TransFrac{\alpha}$ for $\alpha$ sufficiently large.
\end{proof}

\begin{example}[The $\omega$-Fractal]
\label{ex:omega-fractal}
The fractal $\TransFrac{\omega}$ is constructed as:
\[
    \TransFrac{\omega} = \varinjlim_{n < \omega} \TransFrac{n}
\]
This corresponds to the classical infinite iteration of an IFS.
If the base IFS has contraction ratio $r < 1$ with $N$ maps:
\[
    \HausdorffDim(\TransFrac{\omega}) = \frac{\log N}{\log(1/r)}
\]
\end{example}

%% ============================================================================
%% CHAPTER 2: MULTI-INDEXED FRACTAL SYSTEMS
%% ============================================================================

\chapter{Multi-Indexed Fractal Systems}
\label{ch:multi-indexed}

\begin{tcolorbox}[colback=blue!5!white,colframe=blue!75!black,title=\vnewfour]
This chapter generalizes from single ordinal indexing to multi-indexed fractals
$\MultiFrac{\alpha}{\beta}$, enabling richer structural relationships.
\end{tcolorbox}

\section{Definition of $\MultiFrac{\alpha}{\beta}$}

\begin{definition}[Multi-Indexed Fractal]
\label{def:multi-frac}
A \emph{bi-indexed fractal} is a family $\{\MultiFrac{\alpha}{\beta}\}_{(\alpha,\beta) \in \Ord \times \Ord}$
satisfying:
\begin{enumerate}[label=(MF\arabic*)]
    \item \textbf{Vertical Refinement}: For fixed $\beta$, the family
    $\{\MultiFrac{\alpha}{\beta}\}_{\alpha \in \Ord}$ is a transfinite fractal.

    \item \textbf{Horizontal Refinement}: For fixed $\alpha$, the family
    $\{\MultiFrac{\alpha}{\beta}\}_{\beta \in \Ord}$ is a transfinite fractal.

    \item \textbf{Commutation}: The refinement diagram commutes for all
    $\alpha \leq \alpha'$ and $\beta \leq \beta'$.
\end{enumerate}
\end{definition}

\begin{definition}[General Multi-Index]
\label{def:general-multi}
For an index set $I$, an \emph{$I$-indexed fractal} is a functor:
\[
    \Phi^{(\bullet)}: \Ord^I \longrightarrow \FractalFunctor
\]
where $\Ord^I$ is the product category with componentwise ordering.
\end{definition}

\begin{theorem}[Multi-Index Existence]
\label{thm:multi-existence}
Given any collection of transfinite fractals $\{\TransFrac{\alpha}_i\}_{i \in I}$,
there exists a universal $I$-indexed fractal $\Phi^{(\bullet)}$ such that
the $i$-th slice recovers $\TransFrac{\alpha}_i$.
\end{theorem}

\section{Product and Tensor Fractals}

\begin{definition}[Product Fractal]
\label{def:product-frac}
The \emph{product} of fractals $\TransFrac{\alpha}$ and $\TransFrac{\beta}$ is:
\[
    \TransFrac{\alpha} \times_\Phi \TransFrac{\beta} \;:=\;
    \bigl(\TransFrac{\alpha} \times \TransFrac{\beta},\,
    \sigma_\alpha \times \sigma_\beta,\, d_\times\bigr)
\]
where $d_\times((x,y),(x',y')) := \max\{d_\alpha(x,x'), d_\beta(y,y')\}$.
\end{definition}

\begin{proposition}[Product Dimension]
\label{prop:product-dim}
\[
    \HausdorffDim(\TransFrac{\alpha} \times_\Phi \TransFrac{\beta}) \;=\;
    \HausdorffDim(\TransFrac{\alpha}) + \HausdorffDim(\TransFrac{\beta})
\]
\end{proposition}

\begin{definition}[Tensor Fractal]
\label{def:tensor-frac}
The \emph{tensor product} of fractals, denoted $\TransFrac{\alpha} \otimes_\Phi \TransFrac{\beta}$,
is the universal fractal equipped with a bimorphism satisfying the standard
tensor universal property.
\end{definition}

\begin{theorem}[Tensor-Product Adjunction]
\label{thm:tensor-adjunction}
The tensor product $\otimes_\Phi$ is left adjoint to the internal hom
$[\![-,-]\!]_\Phi$ in $\FractalFunctor$:
\[
    \mathrm{Hom}_{\FractalFunctor}(\TransFrac{\alpha} \otimes_\Phi \TransFrac{\beta},\, \Psi)
    \;\cong\;
    \mathrm{Hom}_{\FractalFunctor}(\TransFrac{\alpha},\, [\![\TransFrac{\beta}, \Psi]\!]_\Phi)
\]
\end{theorem}

\section{Fractal Lattice Structures}

\begin{definition}[Fractal Lattice]
\label{def:frac-lattice}
The \emph{fractal lattice} $\mathcal{L}_\Phi$ is the partially ordered set
of all fractals under the refinement ordering:
\[
    \TransFrac{\alpha} \leq_{\mathcal{L}_\Phi} \TransFrac{\beta}
    \quad:\Longleftrightarrow\quad
    \TransFrac{\alpha} \sqsubseteq_{\mathrm{ref}} \TransFrac{\beta}
\]
\end{definition}

\begin{theorem}[Lattice Structure]
\label{thm:lattice-structure}
$(\mathcal{L}_\Phi, \leq_{\mathcal{L}_\Phi})$ is a complete lattice with:
\begin{enumerate}[label=(\alph*)]
    \item \textbf{Meet}: $\bigwedge_i \Phi_i = \bigcap_i \Phi_i$,
    \item \textbf{Join}: $\bigvee_i \Phi_i = \overline{\bigcup_i \Phi_i}$,
    \item \textbf{Bottom}: $\bot = \varnothing$,
    \item \textbf{Top}: $\top = \Ideas$.
\end{enumerate}
\end{theorem}

%% ============================================================================
%% CHAPTER 3: FRACTAL SPECTRUM THEORY
%% ============================================================================

\chapter{Fractal Spectrum Theory}
\label{ch:spectrum}

\begin{tcolorbox}[colback=blue!5!white,colframe=blue!75!black,title=\vnewfour]
This chapter develops spectral methods for fractal operators, including
eigenvalue problems that characterize self-similarity at each level.
\end{tcolorbox}

\section{The Fractal Spectrum $\FractalSpectrum$}

\begin{definition}[Scaling Operator]
\label{def:scaling-op}
For a fractal $\Phi$ with self-similarity map $\sigma: \Phi \to \Phi$,
the \emph{scaling operator} is the induced map on $L^2(\Phi, \mu)$:
\[
    \mathcal{S}_\sigma: L^2(\Phi, \mu) \longrightarrow L^2(\Phi, \mu), \qquad
    (\mathcal{S}_\sigma f)(x) := f(\sigma(x))
\]
where $\mu$ is the natural Hausdorff measure on $\Phi$.
\end{definition}

\begin{definition}[Fractal Spectrum]
\label{def:frac-spectrum}
The \emph{fractal spectrum} of $\Phi$, denoted $\FractalSpectrum$, is the
spectrum of the scaling operator:
\[
    \FractalSpectrum \;:=\; \sigma(\mathcal{S}_\sigma) \;=\;
    \{\lambda \in \mathbb{C} \mid \mathcal{S}_\sigma - \lambda I
    \text{ is not invertible}\}
\]
\end{definition}

\begin{theorem}[Spectral Structure]
\label{thm:spectral-structure}
For a compact fractal $\Phi$ with contraction ratio $r < 1$:
\begin{enumerate}[label=(\alph*)]
    \item $\FractalSpectrum \subseteq \overline{\mathbb{D}}$ (closed unit disk),
    \item $\FractalSpectrum$ is compact,
    \item $1 \in \FractalSpectrum$ with eigenfunction the constant function,
    \item If $\sigma$ has $N$ branches with equal contraction $r$, then
          $r^{-s} \in \FractalSpectrum$ where $s = \HausdorffDim(\Phi)$.
\end{enumerate}
\end{theorem}

\section{Spectral Decomposition of Fractal Operators}

\begin{theorem}[Spectral Decomposition]
\label{thm:spectral-decomp}
Let $\mathcal{S}_\sigma$ be the scaling operator on a fractal $\Phi$. There
exists a spectral measure $E$ on $\FractalSpectrum$ such that:
\[
    \mathcal{S}_\sigma = \int_{\FractalSpectrum} \lambda \, dE(\lambda)
\]
\end{theorem}

\begin{definition}[Multi-Scale Spectral Operator]
\label{def:multi-scale-spectral}
For a transfinite fractal $\{\TransFrac{\alpha}\}_{\alpha \in \Ord}$,
the \emph{multi-scale spectral operator} is:
\[
    \mathcal{S}^{\mathrm{multi}} := \bigoplus_{\alpha \in \Ord} \mathcal{S}_{\sigma_\alpha}
\]
acting on the direct sum Hilbert space.
\end{definition}

\begin{theorem}[Multi-Scale Spectrum]
\label{thm:multi-scale-spectrum}
\[
    \sigma(\mathcal{S}^{\mathrm{multi}}) = \overline{\bigcup_{\alpha \in \Ord}
    \sigma(\mathcal{S}_{\sigma_\alpha})}
\]
\end{theorem}

\section{Eigenvalue Problems for Self-Similarity}

\begin{definition}[Fractal Eigenvalue Problem]
\label{def:frac-eigenvalue}
The \emph{fractal eigenvalue problem} is to find $\lambda_\Phi \in \mathbb{C}$
and nonzero $f \in L^2(\Phi, \mu)$ such that:
\[
    \mathcal{S}_\sigma f = \lambda_\Phi f
\]
\end{definition}

\begin{theorem}[Eigenvalue Characterization]
\label{thm:eigenvalue-char}
For a self-similar fractal with IFS $\{w_1, \ldots, w_N\}$ each with
contraction ratio $r$:
\begin{enumerate}[label=(\alph*)]
    \item The eigenvalues are of the form $\lambda_\Phi = r^k$ for $k \in \mathbb{Z}_{\geq 0}$.
    \item The multiplicity of $\lambda_\Phi = r^k$ is bounded by $N^k$.
    \item The eigenfunctions for $\lambda_\Phi = r^k$ are polynomials of
          degree at most $k$ in suitable fractal coordinates.
\end{enumerate}
\end{theorem}

\begin{definition}[Spectral Zeta Function]
\label{def:spectral-zeta}
The \emph{spectral zeta function} of a fractal $\Phi$ is:
\[
    \zeta_\Phi(s) := \sum_{\lambda_\Phi \in \sigma_p(\mathcal{S}_\sigma)}
    \lambda_\Phi^{-s} = \sum_{k=0}^\infty m_k \cdot r^{-ks}
\]
where $m_k$ is the multiplicity of eigenvalue $r^k$.
\end{definition}

\begin{theorem}[Zeta-Dimension Relation]
\label{thm:zeta-dim}
The spectral zeta function has a pole at $s = \HausdorffDim(\Phi)$:
\[
    \HausdorffDim(\Phi) = \inf\{s > 0 \mid \zeta_\Phi(s) < \infty\}
\]
\end{theorem}

\begin{example}[Sierpinski Triangle Spectrum]
\label{ex:sierpinski-spectrum}
The Sierpinski triangle has $N = 3$ maps with $r = 1/2$. The spectral zeta function:
\[
    \zeta_{\mathrm{ST}}(s) = \frac{1}{1 - 3 \cdot 2^s}
\]
has a pole at $s = \log_2 3 \approx 1.585$.
\end{example}

%% ============================================================================
%% CHAPTER 4: FUNCTORIAL FRACTAL CONSTRUCTIONS
%% ============================================================================

\chapter{Functorial Fractal Constructions}
\label{ch:functorial}

\begin{tcolorbox}[colback=blue!5!white,colframe=blue!75!black,title=\vnewfour]
This chapter establishes the category $\FractalFunctor$ of fractals and studies
functors, natural transformations, and categorical limits/colimits.
\end{tcolorbox}

\section{The Category $\FractalFunctor$ of Fractals}

\begin{definition}[Category of Fractals]
\label{def:frac-cat}
The category $\FractalFunctor$ consists of:
\begin{itemize}
    \item \textbf{Objects}: Fractals $\Phi = (X, \sigma, d)$ where $X$ is a
    compact metric space, $\sigma: X \to X$ is a contractive self-similarity,
    and $d$ is the metric.

    \item \textbf{Morphisms}: A morphism $f: \Phi \to \Psi$ is a continuous
    map that \emph{intertwines} the self-similarities:
    \[
    \begin{tikzcd}
        X \arrow[r, "f"] \arrow[d, "\sigma_\Phi"'] & Y \arrow[d, "\sigma_\Psi"] \\
        X \arrow[r, "f"'] & Y
    \end{tikzcd}
    \]
\end{itemize}
\end{definition}

\begin{theorem}[Categorical Properties of $\FractalFunctor$]
\label{thm:cat-props}
The category $\FractalFunctor$ is:
\begin{enumerate}[label=(\alph*)]
    \item \textbf{Complete}: All small limits exist.
    \item \textbf{Cocomplete}: All small colimits exist.
    \item \textbf{Cartesian closed}: Internal homs exist.
    \item \textbf{Locally small}: Hom-sets are sets.
\end{enumerate}
\end{theorem}

\section{Fractal Functors Between World Categories}

\begin{definition}[Fractal Functor]
\label{def:frac-functor}
A \emph{fractal functor} $\mathbf{F}: \mathbf{World} \to \FractalFunctor$ assigns:
\begin{itemize}
    \item To each world $w$, a fractal $\mathbf{F}(w)$.
    \item To each accessibility $R: w \to w'$, a fractal morphism
          $\mathbf{F}(R): \mathbf{F}(w) \to \mathbf{F}(w')$.
\end{itemize}
\end{definition}

\begin{example}[Refinement Functor]
\label{ex:refinement-functor}
The \emph{refinement functor}:
\[
    \mathbf{Ref}: \Ord \to \FractalFunctor, \qquad \alpha \mapsto \TransFrac{\alpha}
\]
sends each ordinal to its corresponding ordinal fractal.
\end{example}

\section{Natural Transformations of Fractal Systems}

\begin{definition}[Natural Transformation of Fractals]
\label{def:nat-trans}
Let $\mathbf{F}, G: \mathcal{C} \to \FractalFunctor$ be fractal functors.
A \emph{natural transformation} $\eta: \mathbf{F} \Rightarrow G$ is a
family of fractal morphisms $\eta_X: \mathbf{F}(X) \to G(X)$ satisfying naturality.
\end{definition}

\begin{theorem}[Vertical Composition]
\label{thm:vertical-comp}
Natural transformations compose vertically: if $\eta: \mathbf{F} \Rightarrow G$
and $\epsilon: G \Rightarrow H$, then $\epsilon \circ \eta: \mathbf{F} \Rightarrow H$.
\end{theorem}

\begin{theorem}[Horizontal Composition]
\label{thm:horizontal-comp}
Natural transformations compose horizontally via the interchange law.
\end{theorem}

\section{Colimits and Limits of Fractal Diagrams}

\begin{definition}[Colimit of Fractals]
\label{def:colimit-frac}
The \emph{colimit} of a diagram $D: J \to \FractalFunctor$, denoted $\ColimitFrac_J D$,
is a fractal $\Phi$ together with morphisms $\iota_j: D(j) \to \Phi$ universal
for cocones.
\end{definition}

\begin{theorem}[Colimit Construction]
\label{thm:colimit-construction}
For a directed diagram $D: J \to \FractalFunctor$, the colimit is:
\[
    \ColimitFrac_J D = \left(\varinjlim_{j \in J} X_j,\,
    \varinjlim \sigma_j,\, d_{\varinjlim}\right)
\]
\end{theorem}

\begin{example}[Colimit as Infinite Iteration]
\label{ex:colimit-iteration}
The transfinite fractal $\TransFrac{\omega}$ is the colimit:
\[
    \TransFrac{\omega} = \varinjlim_{n < \omega} \TransFrac{n}
\]
\end{example}

\begin{theorem}[Limit-Colimit Duality]
\label{thm:lim-colim-duality}
For a diagram $D: J \to \FractalFunctor$:
\[
    \mathrm{Hom}_{\FractalFunctor}\left(\ColimitFrac_J D,\, \Psi\right) \cong
    \varprojlim_J \mathrm{Hom}_{\FractalFunctor}(D(-), \Psi)
\]
\end{theorem}

%% ============================================================================
%% CHAPTER 5: INTEGRATION WITH v7.3 FRACTAL THEORY
%% ============================================================================

\chapter{Integration with v7.3 Fractal Theory}
\label{ch:integration}

\begin{tcolorbox}[colback=green!5!white,colframe=green!75!black,title=\textbf{Compatibility}]
This chapter establishes compatibility of v7.4 constructions with v7.3,
ensuring the theory is genuinely additive.
\end{tcolorbox}

\section{Extension Theorems from v7.3 Fractals}

\begin{theorem}[v7.3 Embedding]
\label{thm:v73-embedding}
There is a full and faithful embedding:
\[
    \iota: \FractalFunctor_{\mathrm{v7.3}} \hookrightarrow \FractalFunctor_{\mathrm{v7.4}}
\]
preserving all v7.3 structures.
\end{theorem}

\begin{corollary}[Backward Compatibility]
\label{cor:backward-compat}
Every theorem of v7.3 fractal theory remains valid in v7.4 when restricted
to single-indexed fractals.
\end{corollary}

\section{Compatibility with Hausdorff Dimension}

\begin{theorem}[Dimension Consistency]
\label{thm:dim-consistency}
For any v7.3 fractal $\Phi$ and its v7.4 embedding $\iota(\Phi)$:
\[
    \HausdorffDim(\Phi) = \HausdorffDim(\iota(\Phi))
\]
\end{theorem}

\begin{theorem}[Multi-Index Dimension Formula]
\label{thm:multi-dim-formula}
For a bi-indexed tensor fractal:
\[
    \HausdorffDim(\MultiFrac{\alpha}{\beta}) = \HausdorffDim(\TransFrac{\alpha})
    + \HausdorffDim(\TransFrac{\beta})
\]
\end{theorem}

\section{Connection to Fractal Attractors}

\begin{theorem}[Attractor Functoriality]
\label{thm:attractor-functor}
The attractor assignment defines a functor:
\[
    \FractalAttractor^{(-)}: \Ord^I \to \FractalFunctor
\]
natural with respect to index morphisms.
\end{theorem}

\begin{theorem}[Colimit Attractor]
\label{thm:colimit-attractor}
For a directed system of fractals:
\[
    \FractalAttractor(\ColimitFrac_J D) = \varinjlim_{j \in J} \FractalAttractor(D(j))
\]
\end{theorem}

\begin{corollary}[Transfinite Attractor]
\label{cor:transfinite-attractor}
For a transfinite fractal:
\[
    \FractalAttractor(\Phi^{\mathrm{trans}}) = \varinjlim_{\alpha \in \Ord}
    \FractalAttractor(\TransFrac{\alpha})
\]
\end{corollary}

%% ============================================================================
%% CHAPTER 6: RPM DEPENDENCY STRUCTURES AND WORLD TRANSFORMATIONS
%% ============================================================================

\chapter{RPM Dependency Structures and World Transformations}
\label{ch:rpm-dependencies}

\begin{tcolorbox}[colback=blue!5!white,colframe=blue!75!black,title=\vnewfour,breakable]
This chapter formalizes the RPM (Reason, Power, Manifestation) dependency tensor
across worlds, foundations, and noetic aspects. We develop the functorial descent
structure that governs how noetic content transforms through the ACBE chain
$A \to C \to B \to E$ (or equivalently $A \to B \to C \to D$).
\end{tcolorbox}

\section{The RPM Dependency Tensor}

The RPM monad operates across all four worlds ($A, B, C, D$) and all seven
foundations ($F_1, \ldots, F_7$). The following table characterizes how each
acquisition stage interacts with foundation-world pairs.

\renewcommand{\arraystretch}{1.35}
\begin{longtable}{p{2cm}|p{2cm}|c|p{5.5cm}}
\caption{RPM Dependency Tensor: Acquisition $\times$ Foundation $\times$ World}
\label{tab:rpm-dependency-tensor} \\
\toprule
\textbf{Acquisition}
& \textbf{Foundation $F_m$}
& \textbf{World}
& \textbf{RPM Dependency Effect $\mathcal{R}(A_n, F_m, W)$}
\\
\midrule
\endfirsthead
\multicolumn{4}{c}{\textit{Table \ref{tab:rpm-dependency-tensor} continued from previous page}} \\
\toprule
\textbf{Acquisition}
& \textbf{Foundation}
& \textbf{W}
& \textbf{RPM Dependency Effect}
\\
\midrule
\endhead
\midrule
\multicolumn{4}{r}{\textit{Continued on next page}} \\
\endfoot
\bottomrule
\endlastfoot

%==================== A0 — Pure Desire ====================
A0 (Pure Desire) & $F_1$ (Unity) & A
& \textbf{Strengthened.} Archetypal purification ensures $A0$ is uncontradicted; reduces RPM failure risk at first step. \\

A0 & $F_5$ (Power) & C
& \textbf{Amplified.} Emotional sovereignty clarifies true intent; prevents false-desire RPM loops. \\

A0 & $F_6$ (Wealth) & D
& \textbf{Grounded.} Reformulates desire into materially interpretable goals. \\

\midrule

%==================== D_n — Desire Level ====================
$D_n$ & $F_2$ (Wisdom) & B
& \textbf{Required.} Cognitive structuring resolves ambiguous desire; ensures RPM can proceed to $W_n$. \\

$D_n$ & $F_3$ (Life) & C
& \textbf{Strengthened.} Adds emotional vitality to keep $D_n$ from collapsing under conflict. \\

$D_n$ & $F_7$ (Continuation) & D
& \textbf{Extended.} Converts $D_n$ into long-term desire, enabling sustained RPM evaluation. \\

\midrule

%==================== W_n — Wisdom Level ====================
$W_n$ & $F_1$ (Unity) & A
& \textbf{Stabilized.} Archetypal integration of knowledge removes contradictory interpretations. \\

$W_n$ & $F_4$ (Companionship) & C
& \textbf{Contextualized.} Adds relational information to correct misaligned interpretations. \\

$W_n$ & $F_5$ (Power) & B
& \textbf{Sharpened.} Strategic coherence increases the likelihood that $W_n \gg= P_n$ succeeds. \\

\midrule

%==================== P_n — Power Level ====================
$P_n$ & $F_3$ (Life) & D
& \textbf{Energized.} Converts the power expression into kinetic (actionable) form. \\

$P_n$ & $F_6$ (Wealth) & B
& \textbf{Optimized.} Resource-efficiency constraints rewrite $P_n$ for maximum feasibility. \\

$P_n$ & $F_7$ (Continuation) & C
& \textbf{Sustained.} Emotional drive prevents abandonment of the $P_n$ pathway. \\

\midrule

%==================== FAILURE CONDITIONS ====================
\textbf{Any Stage} & $F_{m,w}$ & C
& \textbf{May Block.} If emotional interpretation contradicts cognitive or causal vectors, RPM halts with \emph{Failure}(A). \\

\textbf{Any Stage} & $F_{m,w}$ & D
& \textbf{May Collapse.} If material feasibility is violated, RPM cannot enter $P_n$. \\

\textbf{Any Stage} & $F_{m,w}$ & A
& \textbf{Overrides.} Causal archetype forces re-evaluation of prior stages; RPM backtracks safely. \\

\end{longtable}

\section{Sub-Foundation Effects on Transfinite Fractals}

The sub-foundations $F_{mw}$ (foundation $m$ in world $w$) induce specific
effects on transfinite fractal evaluation. The following table characterizes
behavior at ordinal levels $\omega$, $\omega^2$, and $\omega^\omega$.

\renewcommand{\arraystretch}{1.35}
\begin{longtable}{p{2.2cm}|c|p{1.5cm}|p{5.5cm}}
\caption{Sub-Foundation Effects on Transfinite Fractals $(\omega, \omega^2, \omega^\omega)$}
\label{tab:subfoundation-transfinite-fractals} \\
\toprule
\textbf{Sub-Found.\ $F_{mw}$}
& \textbf{W}
& \textbf{Fractal}
& \textbf{Effect on Transfinite Evaluation}
\\
\midrule
\endfirsthead
\multicolumn{4}{c}{\textit{Table \ref{tab:subfoundation-transfinite-fractals} continued from previous page}} \\
\toprule
\textbf{Sub-Found.\ $F_{mw}$}
& \textbf{W}
& \textbf{Fractal}
& \textbf{Effect on Transfinite Evaluation}
\\
\midrule
\endhead
\midrule
\multicolumn{4}{r}{\textit{Continued on next page}} \\
\endfoot
\bottomrule
\endlastfoot

%==================== UNITY ====================
$F_{1a}$ (Unity--A) & A & $\omega$
& Forces stabilization of the $\omega$-chain at the archetypal fixed point:
  $\sem{\TransFrac{\omega}} = \operatorname{lfp}(\nu_k)$ under world-A constraints. \\

$F_{1a}$ & A & $\omega^2$
& Collapses each $\omega$-block to its archetypal representative, producing:
  $\sem{\TransFrac{\omega^2}} = \bigcup_{i < \omega} \operatorname{lfp}(\Phi_i)$. \\

$F_{1a}$ & A & $\omega^\omega$
& Forces the entire fractal tower toward a single global archetypal attractor;
  all limit ordinals map to the same causal fixed point. \\

\midrule

$F_{1d}$ (Unity--D) & D & $\omega$
& Enforces physical-feasibility truncation:
  the $\omega$-sequence stabilizes early due to material-world grounding. \\

$F_{1d}$ & D & $\omega^2$
& Each $\omega$-block is truncated at the earliest physically admissible point. \\

\midrule

%==================== WISDOM ====================
$F_{2b}$ (Wisdom--B) & B & $\omega$
& Introduces cognitive normalization between each iterate:
  $\nu_k^n \mapsto \mathrm{nf}(\nu_k^n)$, accelerating convergence. \\

$F_{2b}$ & B & $\omega^2$
& Forces each $\omega$-block to converge before the next block begins;
  resembles ordinal-indexed Newton iteration. \\

$F_{2b}$ & B & $\omega^\omega$
& Produces a stratified hierarchy of cognitive consistencies;
  each tower level is normalized before transfinite ascent. \\

\midrule

%==================== LIFE ====================
$F_{3c}$ (Life--C) & C & $\omega$
& Emotional amplification introduces oscillatory behavior;
  $\omega$-limit may be periodic rather than fixed. \\

$F_{3c}$ & C & $\omega^2$
& Creates nested oscillation layers;
  results in a periodic-within-periodic convergence pattern. \\

$F_{3c}$ & C & $\omega^\omega$
& Constructs an infinitely nested emotional resonance stack;
  may block convergence unless stabilized by $F_{1a}$ or $F_{2b}$. \\

\midrule

%==================== COMPANIONSHIP ====================
$F_{4b}$ (Comp.--B) & B & $\omega$
& Introduces theory-of-mind correction to each iterate;
  convergence shifts toward relational equilibrium. \\

$F_{4b}$ & B & $\omega^2$
& Forces each block to account for relational constraints before the next begins. \\

\midrule

%==================== POWER ====================
$F_{5d}$ (Power--D) & D & $\omega$
& Converts the sequence into physical action profiles;
  may truncate non-implementable iterates. \\

$F_{5d}$ & D & $\omega^2$
& Each $\omega$-block maps to a full material action sequence;
  higher blocks represent higher-level plans. \\

$F_{5d}$ & D & $\omega^\omega$
& Generates an infinitely deep hierarchy of physical action schemas---
  a complete ``power-tower'' of embodiment. \\

\midrule

%==================== WEALTH ====================
$F_{6b}$ (Wealth--B) & B & $\omega$
& Convergence weighted by value/utility function;
  limit ordinal selects maximal-utility fixed point. \\

$F_{6b}$ & B & $\omega^2$
& Produces a multi-layer optimization hierarchy;
  each block performs a new level of resource rationalization. \\

$F_{6b}$ & B & $\omega^\omega$
& Forms an infinite optimization tower analogous to
  a transfinite Bellman equation. \\

\midrule

%==================== CONTINUATION ====================
$F_{7a}$ (Cont.--A) & A & $\omega$
& Ensures $\omega$ iteration never collapses prematurely;
  extends sequence to full limit ordinal. \\

$F_{7a}$ & A & $\omega^2$
& Guarantees progression through each $\omega$-block;
  prevents early stabilization. \\

$F_{7a}$ & A & $\omega^\omega$
& Forces the entire fractal tower to express continuation;
  produces genuinely transfinite evolution without early convergence. \\

\end{longtable}

\section{Sub-Foundations as Natural Transformations}

The sub-foundations can be understood category-theoretically as natural
transformations on the category $\Noetica$. Each $F_{mw}$ induces a functor
with specific naturality conditions encoding noetic semantics.

\renewcommand{\arraystretch}{1.35}
\begin{longtable}{p{2cm}|p{2cm}|p{2.5cm}|p{4.5cm}}
\caption{Sub-Foundations as Natural Transformations on the Category Noetica}
\label{tab:subfoundation-natural-transformations} \\
\toprule
\textbf{Sub-Found.\ $F_{mw}$}
& \textbf{Functor}
& \textbf{Naturality}
& \textbf{Noetic Interpretation}
\\
\midrule
\endfirsthead
\multicolumn{4}{c}{\textit{Table \ref{tab:subfoundation-natural-transformations} continued from previous page}} \\
\toprule
\textbf{Sub-Found.\ $F_{mw}$}
& \textbf{Functor}
& \textbf{Naturality}
& \textbf{Noetic Interpretation}
\\
\midrule
\endhead
\midrule
\multicolumn{4}{r}{\textit{Continued on next page}} \\
\endfoot
\bottomrule
\endlastfoot

%==================== UNITY ====================
$F_{1a}$ (Unity--A)
& $\mathcal{U}_A : \Noetica \to \Noetica$
& $\mathcal{U}_A(\nu_k) \circ \eta_x = \eta_y \circ \nu_k$
& Unity lifts every morphism into the archetypal fiber.
  Ensures \emph{global coherence} of all world-A transformations.
  The naturality square asserts:
  ``Noetic operations commute with unity-projection.'' \\

$F_{1d}$ (Unity--D)
& $\mathcal{U}_D$
& Same as above, restricted to $\mathcal{I}_D$.
& Forces grounding: all naturality squares collapse onto the physical-world image.
\\

\midrule

%==================== WISDOM ====================
$F_{2b}$ (Wisdom--B)
& $\mathcal{W}_B$
& $\mathcal{W}_B(\nu_k)(\eta_x) = \eta_{\nu_k(x)}$
& Wisdom enforces \emph{cognitive normalization}.
  Naturality = all noetic operations preserve structural consistency across rational models.
\\

\midrule

%==================== LIFE ====================
$F_{3c}$ (Life--C)
& $\mathcal{L}_C$
& $\eta_{\nu_k(x)} \sqsubseteq \mathcal{L}_C(\nu_k)(\eta_x)$
& Naturality is \emph{inequality} not equality:
  Life introduces \textbf{emotional amplification}, producing monotone growth in the dcpo.
\\

\midrule

%==================== COMPANIONSHIP ====================
$F_{4b}$ (Comp.--B)
& $\mathcal{C}_B$
& $\mathcal{C}_B(\nu_k)(\eta_x)$ reflects relational context
& Companionship modifies morphisms via \emph{relational liftings}.
  Naturality asserts: relational corrections distribute over all noetic transformations.
\\

\midrule

%==================== POWER ====================
$F_{5d}$ (Power--D)
& $\mathcal{P}_D$
& $\eta_y = \mathcal{P}_D(\nu_k)(\eta_x)$ iff action is implementable
& Power enforces \emph{feasibility constraints}.
  The naturality square commutes only when the resulting noetic action is physically executable.
\\

\midrule

%==================== WEALTH ====================
$F_{6b}$ (Wealth--B)
& $\mathcal{M}_B$ (``Material Optimization Functor'')
& $\eta_{\nu_k(x)} = \operatorname*{argmax}(\mathcal{M}_B(\nu_k)(\eta_x))$
& Wealth transforms each morphism into its \textbf{utility-optimized version}.
  Naturality = noetic operations preserve utility maximization under $F_{6b}$.
\\

\midrule

%==================== CONTINUATION ====================
$F_{7a}$ (Continuation--A)
& $\mathcal{K}_A$ (``Transfinite Continuation Functor'')
& $\eta_x \mapsto \bigsqcup_{\alpha < \lambda} \eta_{\TransFrac{\alpha}(x)}$
& Natural transformation that enforces \textbf{non-premature termination}.
  The naturality square expresses:
  ``Every noetic transformation lifts to the transfinite continuation tower.''
\\

\end{longtable}

\section{The RPM Dependency Tensor $\mathbb{T}_{RPM}$}

We now formalize the full tensor structure that captures RPM dependencies
across acquisition nodes, worlds, foundations, and noetic aspect profiles.

\renewcommand{\arraystretch}{1.28}
\begin{longtable}{p{2.2cm}|p{2cm}|p{1.8cm}|p{5cm}}
\caption{The RPM Dependency Tensor $\mathbb{T}_{RPM}$ Across Worlds, Foundations, and Aspects}
\label{tab:RPM-tensor} \\
\toprule
\textbf{Acquisition}
& \textbf{World $W$}
& \textbf{Found.\ $F_m$}
& \textbf{Aspect Profile (Preconditions)}
\\
\midrule
\endfirsthead
\multicolumn{4}{c}{\textit{Table \ref{tab:RPM-tensor} continued from previous page}} \\
\toprule
\textbf{Acquisition}
& \textbf{World $W$}
& \textbf{Found.\ $F_m$}
& \textbf{Aspect Profile (Preconditions)}
\\
\midrule
\endhead
\midrule
\multicolumn{4}{r}{\textit{Continued on next page}} \\
\endfoot
\bottomrule
\endlastfoot

%==================== A0 ====================
$\mathbf{A0}$ (Primordial Desire)
& All $W \in \{A,B,C,D\}$
& None (pre-foundational)
& Requires the \textbf{Desire Aspect} vector:
\[
\vec{a}_{A0} =
\langle \nu_1 \;\text{(Mind)},\;
        \nu_2 \;\text{(Positive)},\;
        \nu_4 \;\text{(Vibration)} \rangle
\]
Tensor slot:
\[
\mathbb{T}_{RPM}[A0, W, F_m] =
\begin{cases}
1 & \text{if $m=1$ or unscoped} \\
0 & \text{otherwise}
\end{cases}
\]
\\

\midrule

%==================== D_n ====================
$\mathbf{D_n}$ (Desire---Foundation $F_n$)
& $W_d \in \{A,B,C,D\}$
& $F_n$
& Aspect profile:
\[
\vec{a}_{D_n} =
\langle
\nu_1,\;
\nu_2,\;
\nu_5
\rangle
\]
Dependency condition:
\[
\mathbb{T}_{RPM}[D_n, W_d, F_n] = 1
\quad\text{iff}\quad
A0 \Downarrow \text{true under } \nu_1,\nu_2,\nu_5.
\]
All off-diagonal tensor entries are $0$.
\\

\midrule

%==================== W_n ====================
$\mathbf{W_n}$ (Wisdom---Foundation $F_n$)
& $W_w \in \{A,B,C,D\}$
& $F_n$
& Aspect constraints:
\[
\vec{a}_{W_n}
=
\langle
\nu_1,\;
\nu_6,\;
\nu_7
\rangle
\]
Dependency rule:
\[
D_n \Rightarrow W_n
\quad\text{iff}\quad
\mathfrak{R}(D_n) \text{ is non-failure.}
\]
Tensor slice:
\[
\mathbb{T}_{RPM}[W_n, W_w, F_n] =
\mathbf{ACBE}(W_n)
\]
indicating dependence on world-transformation functor action.
\\

\midrule

%==================== P_n ====================
$\mathbf{P_n}$ (Power---Foundation $F_n$)
& $W_p \in \{A,B,C,D\}$
& $F_n$
& Aspect profile:
\[
\vec{a}_{P_n} =
\langle
\nu_1,\;
\nu_3,\;
\nu_6
\rangle
\]
Dependency condition:
\[
W_n \Rightarrow P_n
\quad\text{iff}\quad
\exists \lambda.\;
\sem{\TransFrac{\lambda}}(W_n)
\text{ stabilizes.}
\]
Tensor expansion:
\[
\mathbb{T}_{RPM}[P_n, W_p, F_n] =
\mathcal{K}_A(W_n) \otimes \mathcal{P}_D(W_p)
\]
indicating interaction of continuation and feasibility constraints.
\\

\end{longtable}

\section{ACBE Functorial Descent}

The ACBE chain ($A \to C \to B \to E$, or in world notation $A \to B \to C \to D$)
represents the descent of noetic content from archetypal to physical worlds.
Each step is a functor with naturality conditions ensuring noetic operations
are preserved.

\renewcommand{\arraystretch}{1.32}
\small
\begin{longtable}{p{1.2cm}|p{2.2cm}|p{2.8cm}|p{6cm}}
\caption{ACBE Functorial Descent Cube: Natural Transformation Structure Across Worlds}
\label{tab:ACBE-descent-cube} \\
\toprule
\textbf{Worlds}
& \textbf{Functor $\mathcal{F}_W$}
& \textbf{Descent $\delta$}
& \textbf{Naturality Condition}
\\
\midrule
\endfirsthead
\multicolumn{4}{c}{\textit{Table \ref{tab:ACBE-descent-cube} continued from previous page}} \\
\toprule
\textbf{Worlds}
& \textbf{Functor $\mathcal{F}_W$}
& \textbf{Descent $\delta$}
& \textbf{Naturality Condition}
\\
\midrule
\endhead
\midrule
\multicolumn{4}{r}{\textit{Continued on next page}} \\
\endfoot
\bottomrule
\endlastfoot

%------------------ A -> B ------------------
$\mathbf{A \to B}$
& $\mathcal{F}_A(X) = \sem{X}_A$
& $\delta_{A\to B} : \sem{X}_A \to \sem{X}_B$
& For all $x \in \sem{X}_A$:
$\delta_{A\to B}(\nu_k(x)) = \nu_k( \delta_{A\to B}(x) )$.
This encodes the \textbf{Causal $\to$ Mental projection} of Ideas.
\\
\midrule

%------------------ B -> C ------------------
$\mathbf{B \to C}$
& $\mathcal{F}_B(X) = \sem{X}_B$
& $\delta_{B\to C} : \sem{X}_B \to \sem{X}_C$
& Naturality square:
$\begin{tikzcd}[cramped,sep=small]
\sem{X}_B \arrow[r,"\nu_k"] \arrow[d,"\delta_{B\to C}"']
& \sem{X}_B \arrow[d,"\delta_{B\to C}"] \\
\sem{X}_C \arrow[r,"\nu_k"']
& \sem{X}_C
\end{tikzcd}$
Emotion-level compression preserves Noetic action.
\\
\midrule

%------------------ C -> D ------------------
$\mathbf{C \to D}$
& $\mathcal{F}_C(X) = \sem{X}_C$
& $\delta_{C\to D} : \sem{X}_C \to \sem{X}_D$
& Concrete semantics:
$\delta_{C\to D}(x) = \text{``Physical projection of emotional signature''}$
and naturality:
$\delta_{C\to D}(\nu_k(x)) = \nu_k(\delta_{C\to D}(x))$
when $\nu_k$ is physically realizable (i.e., $k \neq 8,9$).
\\
\midrule

%------------------ A -> C ------------------
$\mathbf{A \to C}$
& Composite functor $\mathcal{F}_C \circ \mathcal{F}_B$
& $\delta_{A\to C} = \delta_{B\to C} \circ \delta_{A\to B}$
& Composite naturality:
$\delta_{A\to C}(\nu_k(x)) = \nu_k(\delta_{A\to C}(x))$.
This is the \textbf{2-step descent} of spiritual ideas into emotional form.
\\
\midrule

%------------------ B -> D ------------------
$\mathbf{B \to D}$
& Composite functor $\mathcal{F}_D \circ \mathcal{F}_C$
& $\delta_{B\to D} = \delta_{C\to D} \circ \delta_{B\to C}$
& Naturality preserves computational closure under RPM-evaluable expressions.
\\
\midrule

%------------------ A -> D ------------------
$\mathbf{A \to D}$
& $\mathcal{F}_D(X) = \sem{X}_D$
& $\delta_{A\to D} = \delta_{C\to D} \circ \delta_{B\to C} \circ \delta_{A\to B}$
& Global naturality:
$\delta_{A\to D}(\nu_k(x)) = \nu_k(\delta_{A\to D}(x))$
encodes the \textbf{full ACBE chain}:
$A \xrightarrow{F_B} B \xrightarrow{F_C} C \xrightarrow{F_D} D$.
\\

\end{longtable}

\begin{theorem}[ACBE Descent Coherence]
\label{thm:acbe-coherence}
The ACBE descent morphisms satisfy the coherence condition:
\[
\delta_{A\to D} = \delta_{C\to D} \circ \delta_{B\to C} \circ \delta_{A\to B}
= \delta_{B\to D} \circ \delta_{A\to B}
= \delta_{C\to D} \circ \delta_{A\to C}
\]
All paths through the descent cube yield the same composite morphism.
\end{theorem}

\begin{proof}
By functoriality of composition and the commutativity of the naturality squares
established in each row of Table~\ref{tab:ACBE-descent-cube}.
\end{proof}

\begin{corollary}[RPM-ACBE Integration]
\label{cor:rpm-acbe}
The RPM monad $\mathfrak{R}$ is compatible with ACBE descent:
\[
\delta_{W\to W'}(\mathfrak{R}(x)) = \mathfrak{R}(\delta_{W\to W'}(x))
\]
for all valid acquisition nodes and world pairs.
\end{corollary}

%% ============================================================================
%% CHAPTER 7: NOETIC OPERATOR ALGEBRA AND FRACTAL CLASSIFICATION
%% ============================================================================

\chapter{Noetic Operator Algebra and Fractal Classification}
\label{ch:noetic-algebra}

\begin{tcolorbox}[colback=blue!5!white,colframe=blue!75!black,title=\vnewfour,breakable]
This chapter provides comprehensive classification tables for noetic operators,
their world interactions, duality structures, fractal normal forms, and
transfinite convergence behavior.
\end{tcolorbox}

\section{Noetic--World Interaction Matrix}

The following table characterizes how each noetic operator $\nu_k$ behaves
across the four worlds, including stability properties and descent compatibility.

\renewcommand{\arraystretch}{1.32}
\small
\begin{longtable}{p{1.5cm}|c|c|c|c|p{4.5cm}}
\caption{Noetic--World Interaction Matrix: Cross-World Behavior of Each Operator $\nu_k$}
\label{tab:noetic-world-matrix} \\
\toprule
\textbf{Noetic $\nu_k$}
& \textbf{A}
& \textbf{B}
& \textbf{C}
& \textbf{D}
& \textbf{Notes (Stability, Descent, RPM)}
\\
\midrule
\endfirsthead
\multicolumn{6}{c}{\textit{Table \ref{tab:noetic-world-matrix} continued from previous page}} \\
\toprule
\textbf{Noetic $\nu_k$}
& \textbf{A}
& \textbf{B}
& \textbf{C}
& \textbf{D}
& \textbf{Notes}
\\
\midrule
\endhead
\midrule
\multicolumn{6}{r}{\textit{Continued on next page}} \\
\endfoot
\bottomrule
\endlastfoot

$\nu_0$ (Identity)
& $A \to A$
& $B \to B$
& $C \to C$
& $D \to D$
& World-stable; commutes with all descent morphisms.
Trivial fixed point: $x = \nu_0(x)$.
\\
\midrule

$\nu_1$ (Mind)
& $A \to A$
& $B \to B$
& $C \to B$
& $D \to B$
& Pulls all worlds toward Mental layer.
Stabilizes on $B$; forces RPM cognitive-check stage.
\\
\midrule

$\nu_2$ (Positive)
& $A \to A$
& $B \to B$
& $C \to C$
& $D \to C$
& Upward emotionalization.
At physical level induces attraction toward emotional meaning.
Cancels with $\nu_3$ under duality.
\\
\midrule

$\nu_3$ (Negative)
& $A \to A$
& $B \to B$
& $C \to C$
& $D \to C$
& Counterpart of $\nu_2$.
Drives separation and analytic differentiation.
RPM interprets $\nu_3$ as ``prerequisite obstruction.''
\\
\midrule

$\nu_4$ (Vibration)
& $A \to A$
& $B \to B$
& $C \to C$
& $D \to D$
& Fully world-stable operator.
Generates eigenfractals; key for MVR stabilizing sequences.
Under descent: $\delta(\nu_4(x)) = \nu_4(\delta(x))$.
\\
\midrule

$\nu_5$ (Female)
& $A \to A$
& $B \to C$
& $C \to C$
& $D \to C$
& Lowers world-level via internalization.
Stabilizes on Emotional world.
Dual to $\nu_6$; cancels in composition.
\\
\midrule

$\nu_6$ (Male)
& $A \to B$
& $B \to D$
& $C \to D$
& $D \to D$
& Pushes Ideas downward into embodiment.
Becomes a physicalizing operator under descent.
RPM treats $\nu_6$ as ``intent activation.''
\\
\midrule

$\nu_7$ (Rhythm)
& $A \to A$
& $B \to B$
& $C \to C$
& $D \to D$
& Fully world-stable.
Key for periodic evaluation in fractals and RPM stepping.
Generates limit cycles in emotional and physical domains.
\\
\midrule

$\nu_8$ (Above)
& $A \to A$
& $B \to A$
& $C \to A$
& $D \to A$
& Strong lifting operator.
Inserts spiritual interpretation into all lower-world Ideas.
Non-physical Noetic; does not preserve descent unless trivial.
\\
\midrule

$\nu_9$ (Below)
& $A \to D$
& $B \to D$
& $C \to D$
& $D \to D$
& Strong grounding operator.
Projects all Ideas into physical realizability.
Cancels $\nu_8$ under duality.
Required for RPM ``Power'' checks.
\\

\end{longtable}
\normalsize

\section{Noetic Duality and Mirror Algebra}

\renewcommand{\arraystretch}{1.28}
\small
\begin{longtable}{p{1.3cm}|c|p{1.5cm}|p{1.8cm}|p{2cm}|p{4.5cm}}
\caption{Noetic Duality \& Mirror Algebra: Symmetry Classes, Dual Pairs, Collapse Rules}
\label{tab:noetic-duality-algebra} \\
\toprule
\textbf{$\nu_k$}
& \textbf{Dual}
& \textbf{Mirror}
& \textbf{Fixed Pt}
& \textbf{Collapse}
& \textbf{Algebraic / Semantic Notes}
\\
\midrule
\endfirsthead
\multicolumn{6}{c}{\textit{Table \ref{tab:noetic-duality-algebra} continued from previous page}} \\
\toprule
\textbf{$\nu_k$}
& \textbf{Dual}
& \textbf{Mirror}
& \textbf{Fixed Pt}
& \textbf{Collapse}
& \textbf{Notes}
\\
\midrule
\endhead
\midrule
\multicolumn{6}{r}{\textit{Continued on next page}} \\
\endfoot
\bottomrule
\endlastfoot

$\nu_0$
& $\nu_0$
& All
& $x = \nu_0(x)$
& None
& Identity morphism in $\Noetica$. Neutral element for all compositions.
\\
\midrule

$\nu_1$ (Mind)
& none
& $\mathcal{M}_1$
& $A1$
& $\nu_1^2 = \nu_1$
& Projects all worlds to Mental layer. Unique stabilizer of analytic states.
\\
\midrule

$\nu_2$ (Pos)
& $\nu_3$
& $\mathcal{M}_2$
& positive polarity
& $\nu_2 \circ \nu_3 = \nu_0$
& Attraction operator; counterpart to $\nu_3$. Dual cancellation central.
\\
\midrule

$\nu_3$ (Neg)
& $\nu_2$
& $\mathcal{M}_2$
& negative polarity
& $\nu_3 \circ \nu_2 = \nu_0$
& Differentiation operator; introduces contrast. Causes RPM-cascade failures.
\\
\midrule

$\nu_4$ (Vib)
& none
& $\mathcal{M}_4$
& $\nu_4(x) = \lambda x$
& Idempotent
& Generates eigenfractals. Symmetric under world ascent/descent.
\\
\midrule

$\nu_5$ (Fem)
& $\nu_6$
& $\mathcal{M}_5$
& Emotional cores
& $\nu_6 \circ \nu_5 = \nu_0$
& Receptive operator; internalizes structure; dualizes with $\nu_6$ in RPM.
\\
\midrule

$\nu_6$ (Male)
& $\nu_5$
& $\mathcal{M}_5$
& Physical pts
& $\nu_5 \circ \nu_6 = \nu_0$
& Projective operator; embodiment driver; activates Power-stage in RPM.
\\
\midrule

$\nu_7$ (Rhythm)
& none
& $\mathcal{M}_7$
& Limit cycles
& $\nu_7^n \to$ cycle
& Governs periodicity and stabilization in fractal calculus.
\\
\midrule

$\nu_8$ (Above)
& $\nu_9$
& $\mathcal{M}_8$
& $A$-core pts
& $\nu_9 \circ \nu_8 = \nu_0$
& Lifting operator; pushes Ideas up into Spiritual world; non-grounding.
\\
\midrule

$\nu_9$ (Below)
& $\nu_8$
& $\mathcal{M}_8$
& $D$-core pts
& $\nu_8 \circ \nu_9 = \nu_0$
& Grounding operator; collapses abstraction into physical realizability.
\\

\end{longtable}
\normalsize

\section{Fractal Normal Forms}

\renewcommand{\arraystretch}{1.32}
\small
\begin{longtable}{p{2.8cm}|p{1.8cm}|p{2.2cm}|p{1.5cm}|p{4cm}}
\caption{Fractal Normal Forms: Canonical Representatives, Collapse Rules, and Complexity Classes}
\label{tab:fractal-normal-forms} \\
\toprule
\textbf{Fractal Pattern}
& \textbf{Normal Form}
& \textbf{Collapse Rule}
& \textbf{Complexity}
& \textbf{Semantic Interpretation}
\\
\midrule
\endfirsthead
\multicolumn{5}{c}{\textit{Table \ref{tab:fractal-normal-forms} continued from previous page}} \\
\toprule
\textbf{Pattern}
& \textbf{NF}
& \textbf{Collapse}
& \textbf{Complexity}
& \textbf{Interpretation}
\\
\midrule
\endhead
\midrule
\multicolumn{5}{r}{\textit{Continued on next page}} \\
\endfoot
\bottomrule
\endlastfoot

$\langle k \rangle$
& $\langle k \rangle$
& None
& \textsf{Linear}
& Single-step operator; eigenfractals arise when $\nu_k(x)=\lambda x$.
\\
\midrule

$\langle k:k:k\rangle$
& $\langle k \rangle$
& Idempotence: $\nu_k^n = \nu_k$
& \textsf{Stabilizing}
& Iterating a single Noetic converges to its fixed point.
\\
\midrule

$\langle a:b \rangle$ dual pair
& $\langle 0 \rangle$
& $\nu_a\circ\nu_b = \nu_0$
& \textsf{Collapsing}
& Fractal annihilates its polarity; reduces to identity.
\\
\midrule

$\langle i:j \rangle$ non-dual
& $\langle i:j \rangle$
& None
& \textsf{Quadratic}
& Composition of distinct operators yields shape-dependent behavior.
\\
\midrule

$\langle 1:4:7\rangle$ MVR
& $\langle 147 \rangle_{\mathrm{NF}}$
& World-preserving triplet
& \textsf{Stabilizing}
& MVR installs stable Mind--Vibration--Rhythm loop; healing fractal.
\\
\midrule

$\langle 8:1:4:7:9\rangle$ ACBE
& $\langle 81479 \rangle_{\mathrm{NF}}$
& Partial dual collapse
& \textsf{Stab/Proj}
& Models descent from causal to physical worlds.
\\
\midrule

$\langle 2:3:2:3:\dots \rangle$
& $\langle 0 \rangle$
& Periodic collapse
& \textsf{Cyc$\to$Coll}
& Oscillatory sequence with stable annihilation limit.
\\
\midrule

$\langle 5:6:5:6:\dots \rangle$
& $\langle 0 \rangle$
& Gender duality
& \textsf{Cyc$\to$Coll}
& Emotional $\leftrightarrow$ physical oscillation collapses.
\\
\midrule

$\langle 8:9:8:9:\dots \rangle$
& $\langle 0 \rangle$
& $\nu_8\circ\nu_9=\nu_0$
& \textsf{Cyc$\to$Coll}
& Metaphysical expansion/contraction loops.
\\
\midrule

$\langle k_1:\dots:k_n\rangle$
& $\langle\mathrm{NF}\rangle$
& Recursive collapse
& \textsf{Poly/Exp}
& Shape determines convergence; reduces to irreducible subsequence.
\\

\end{longtable}
\normalsize

\section{Transfinite Fractal Limit Behavior}

\renewcommand{\arraystretch}{1.30}
\small
\begin{longtable}{p{2.2cm}|p{1.3cm}|p{1.5cm}|p{1.5cm}|p{1.3cm}|p{4cm}}
\caption{Transfinite Fractal Limit Behavior: Convergence, Stabilization, Fixed-Points}
\label{tab:transfinite-fractal-limits} \\
\toprule
\textbf{Fractal Class}
& \textbf{Stab.\ Ord.}
& \textbf{Conv.\ Class}
& \textbf{Fixed-Pt}
& \textbf{Basin}
& \textbf{Semantic Interpretation}
\\
\midrule
\endfirsthead
\multicolumn{6}{c}{\textit{Table \ref{tab:transfinite-fractal-limits} continued from previous page}} \\
\toprule
\textbf{Class}
& \textbf{Stab.}
& \textbf{Conv.}
& \textbf{FP}
& \textbf{Basin}
& \textbf{Interpretation}
\\
\midrule
\endhead
\midrule
\multicolumn{6}{r}{\textit{Continued on next page}} \\
\endfoot
\bottomrule
\endlastfoot

$\TransFrac{n}$ (finite)
& $< n$
& Strong
& Simple
& Local
& Classical fractal; collapses to NF after finitely many steps.
\\
\midrule

$\TransFrac{\omega}_k$ constant
& $\le 40$
& $\omega$-Conv
& Idempotent
& Global-$k$
& Constant transfinite iteration stabilizes via idempotence.
\\
\midrule

$\langle a:b:\dots \rangle_\omega$ dual
& $2$
& Osc$\to$Coll
& Identity
& Entire
& Dual alternation collapses; limit is $\nu_0$.
\\
\midrule

$\langle i:j:\dots \rangle_\omega$ non-dual
& $< \omega$
& Osc/Asymp
& Cycle/Fixed
& Partitioned
& Non-dual alternations yield limit cycles; emotional/physical resonance.
\\
\midrule

ACBE $\langle 8:1:4:7:9\rangle_\omega$
& $\omega$
& Weak Conv
& Mixed on $D$
& Wide
& Infinite world descent; stabilizes at physical grounding.
\\
\midrule

$\TransFrac{\lambda}$, $\lambda$ limit
& $\lambda$
& Limit-Stable
& DCPO lfp
& Monotone
& Approximants form directed set; limit is lub in $\sem{\Ideas}_\bot$.
\\
\midrule

$\TransFrac{\alpha}$, $\alpha < \varepsilon_0$
& $\alpha^*$
& Hyper-Conv
& Structured
& Layered
& Ordinal decomposition controls convergence speed.
\\
\midrule

$\TransFrac{\alpha}$, $\alpha \ge \varepsilon_0$
& Unbounded
& Pot.\ Divergent
& None
& Reducible
& Exceeds primitive recursive induction; requires compression.
\\
\midrule

Mixed transfinite
& Varies
& Piecewise
& Hybrid
& Partitioned
& Convergence by dominant subsequence; dual segments collapse.
\\

\end{longtable}
\normalsize

\section{Foundation--World--Aspect Compatibility Matrix}

The 28 sub-foundations arise from the $7 \times 4$ combination of foundations
and worlds. The following table characterizes their primary noetic modes and
semantic domain biases.

\renewcommand{\arraystretch}{1.25}
\small
\begin{longtable}{p{1.8cm}|c|c|p{1.5cm}|p{1.5cm}|p{4.5cm}}
\caption{Foundation--World--Aspect Compatibility Matrix (28 Sub-Foundations)}
\label{tab:28-subfoundation-matrix} \\
\toprule
\textbf{Foundation}
& \textbf{W}
& \textbf{ID}
& \textbf{Noetics}
& \textbf{Domain}
& \textbf{Compatibility Notes}
\\
\midrule
\endfirsthead
\multicolumn{6}{c}{\textit{Table \ref{tab:28-subfoundation-matrix} continued from previous page}} \\
\toprule
\textbf{Foundation}
& \textbf{W}
& \textbf{ID}
& \textbf{Noetics}
& \textbf{Domain}
& \textbf{Notes}
\\
\midrule
\endhead
\midrule
\multicolumn{6}{r}{\textit{Continued on next page}} \\
\endfoot
\bottomrule
\endlastfoot

%========== 1: Awareness ==========
\textbf{1 --- Awareness}
& A & 1A &
$\nu_1, \nu_8$ &
$\mathcal{W}$ &
Ideal for initialization; supports high-coherence RPM prerequisites.
\\
& B & 1B &
$\nu_1, \nu_5$ &
$\mathcal{M}$ &
Stabilizes conceptual clarity; increases fractal readability.
\\
& C & 1C &
$\nu_1, \nu_2$ &
$\mathcal{E}$ &
Heightens emotional resolution; lowers semantic noise.
\\
& D & 1D &
$\nu_1, \nu_4$ &
$\mathcal{P}$ &
Grounds Awareness in embodiment; required for full ACBE descent.
\\
\midrule

%========== 2: Positivity ==========
\textbf{2 --- Positivity}
& A & 2A &
$\nu_2, \nu_8$ &
$\mathcal{W}$ &
Spiritual elevation; harmonizes with A-world fractals.
\\
& B & 2B &
$\nu_2, \nu_5$ &
$\mathcal{M}$ &
Enhances RPM optimism gradient; increases fixed-point attraction.
\\
& C & 2C &
$\nu_2, \nu_6$ &
$\mathcal{E}$ &
Emotional uplift; resonance-friendly state for $\omega$-fractals.
\\
& D & 2D &
$\nu_2, \nu_4$ &
$\mathcal{P}$ &
Physical well-being; supports stabilizing fractal limits.
\\
\midrule

%========== 3: Negativity ==========
\textbf{3 --- Negativity}
& A & 3A &
$\nu_3, \nu_8$ &
$\mathcal{W}$ &
Shadow-aspect identification; compatible with RPM error handling.
\\
& B & 3B &
$\nu_3, \nu_6$ &
$\mathcal{M}$ &
Detects cognitive dissonance; essential for debugging.
\\
& C & 3C &
$\nu_3, \nu_7$ &
$\mathcal{E}$ &
Emotional turbulence mapping; precursor to stabilization.
\\
& D & 3D &
$\nu_3, \nu_4$ &
$\mathcal{P}$ &
Maps somatic stress; useful for embodiment diagnostics.
\\
\midrule

%========== 4: Vibration ==========
\textbf{4 --- Vibration}
& A & 4A &
$\nu_4, \nu_8$ &
$\mathcal{W}$ &
Spiritual frequency tuning; prepares world-descent pathways.
\\
& B & 4B &
$\nu_4, \nu_5$ &
$\mathcal{M}$ &
Cognitive modulation; improves fractal signal-to-noise ratio.
\\
& C & 4C &
$\nu_4, \nu_7$ &
$\mathcal{E}$ &
Emotional resonance amplification; interacts with oscillatory fractals.
\\
& D & 4D &
$\nu_4$ &
$\mathcal{P}$ &
Physical frequency and energy-body entrainment.
\\
\midrule

%========== 5: Receptivity / Feminine ==========
\textbf{5 --- Receptive}
& A & 5A &
$\nu_5, \nu_8$ &
$\mathcal{W}$ &
Archetypal Amma; increases higher-world receptivity.
\\
& B & 5B &
$\nu_5, \nu_1$ &
$\mathcal{M}$ &
Subconscious accessibility; ideal for intuitive reasoning.
\\
& C & 5C &
$\nu_5, \nu_2$ &
$\mathcal{E}$ &
Compassion modulation; nurtures emotional convergence.
\\
& D & 5D &
$\nu_5, \nu_4$ &
$\mathcal{P}$ &
Biological receptivity; embodiment of the receptive pole.
\\
\midrule

%========== 6: Projective / Masculine ==========
\textbf{6 --- Projective}
& A & 6A &
$\nu_6, \nu_8$ &
$\mathcal{W}$ &
Archetypal Abba; initiatory force in spiritual domain.
\\
& B & 6B &
$\nu_6, \nu_1$ &
$\mathcal{M}$ &
Analytic projection; ideal for theorem-proving contexts.
\\
& C & 6C &
$\nu_6, \nu_7$ &
$\mathcal{E}$ &
Asserts emotional gradients; high-energy transformations.
\\
& D & 6D &
$\nu_6, \nu_4$ &
$\mathcal{P}$ &
Physical assertion; supports material manifestation.
\\
\midrule

%========== 7: Rhythm ==========
\textbf{7 --- Rhythm}
& A & 7A &
$\nu_7, \nu_8$ &
$\mathcal{W}$ &
Destiny patterning; aligns spiritual rhythms.
\\
& B & 7B &
$\nu_7, \nu_1$ &
$\mathcal{M}$ &
Thought-cycle identification; interacts well with monadic iteration.
\\
& C & 7C &
$\nu_7, \nu_2$ &
$\mathcal{E}$ &
Mood-cycle stabilization; reduces oscillatory divergence.
\\
& D & 7D &
$\nu_7, \nu_4$ &
$\mathcal{P}$ &
Physical cycles (biorhythm, breath, gait); essential for embodiment.
\\

\end{longtable}
\normalsize

\section{Interoperability of Sub-Foundations Under ACBE Functor}

\renewcommand{\arraystretch}{1.25}
\small
\begin{longtable}{c|c|c|c|p{1.5cm}|p{4.5cm}}
\caption{Interoperability of Sub-Foundations Under ACBE Functor}
\label{tab:subfoundation-acbe} \\
\toprule
\textbf{Input}
& \textbf{Functor}
& \textbf{Output}
& \textbf{Pres.}
& \textbf{Noetics}
& \textbf{Semantic Transformation Summary}
\\
\midrule
\endfirsthead
\multicolumn{6}{c}{\textit{Table \ref{tab:subfoundation-acbe} continued from previous page}} \\
\toprule
\textbf{Input}
& \textbf{Functor}
& \textbf{Output}
& \textbf{Pres.}
& \textbf{Noetics}
& \textbf{Summary}
\\
\midrule
\endhead
\midrule
\multicolumn{6}{r}{\textit{Continued on next page}} \\
\endfoot
\bottomrule
\endlastfoot

%----- Awareness (1A–1D)
1A & $F_A$ & 1B & P &
$\nu_1,\nu_8$ &
Spiritual awareness condenses into mental clarity; structure-preserving.
\\
1B & $F_B$ & 1C & P &
$\nu_1,\nu_5$ &
Mental attention descends into emotional recognition; stable transition.
\\
1C & $F_C$ & 1D & W &
$\nu_1,\nu_2$ &
Emotional awareness grounds into physical sensation; meaning weakens.
\\
1D & $F_D$ & 1D & P &
$\nu_1,\nu_4$ &
Identity fixed; terminal stage of descent.
\\
\midrule

%----- Positivity (2A–2D)
2A & $F_A$ & 2B & A &
$\nu_2,\nu_8$ &
Spiritual positivity amplifies as mental optimism; upward bias preserved.
\\
2B & $F_B$ & 2C & P &
$\nu_2,\nu_5$ &
Optimism becomes emotional uplift; structurally preserved.
\\
2C & $F_C$ & 2D & W &
$\nu_2,\nu_6$ &
Emotional positivity becomes physical wellness; slight attenuation.
\\
2D & $F_D$ & 2D & P &
$\nu_2,\nu_4$ &
Fixed point; no further descent.
\\
\midrule

%----- Negativity (3A–3D)
3A & $F_A$ & 3B & A &
$\nu_3,\nu_8$ &
Shadow clarity increases cognitive dissonance detection; amplification.
\\
3B & $F_B$ & 3C & P &
$\nu_3,\nu_6$ &
Mental resistance transforms to emotional turbulence; preserved.
\\
3C & $F_C$ & 3D & A &
$\nu_3,\nu_7$ &
Emotional negativity intensifies into physical stress markers; amplified.
\\
3D & $F_D$ & 3D & P &
$\nu_3,\nu_4$ &
Terminal embodiment; fixed.
\\
\midrule

%----- Vibration (4A–4D)
4A & $F_A$ & 4B & A &
$\nu_4,\nu_8$ &
Spiritual frequency compresses into mental modulation; amplification.
\\
4B & $F_B$ & 4C & P &
$\nu_4,\nu_5$ &
Mental vibration becomes emotional resonance; structurally preserved.
\\
4C & $F_C$ & 4D & P &
$\nu_4,\nu_7$ &
Emotional oscillation grounds to physical frequency; preserved.
\\
4D & $F_D$ & 4D & P &
$\nu_4$ &
Final grounding; remains stable.
\\
\midrule

%----- Receptive (Feminine) (5A–5D)
5A & $F_A$ & 5B & P &
$\nu_5,\nu_8$ &
Archetypal receptivity descends into intuitive cognition; preserved.
\\
5B & $F_B$ & 5C & P &
$\nu_5,\nu_1$ &
Mental openness $\to$ emotional compassion; preserved.
\\
5C & $F_C$ & 5D & W &
$\nu_5,\nu_2$ &
Compassion becomes physical softness; slight semantic attenuation.
\\
5D & $F_D$ & 5D & P &
$\nu_5,\nu_4$ &
Terminal embodiment.
\\
\midrule

%----- Projective (Masculine) (6A–6D)
6A & $F_A$ & 6B & A &
$\nu_6,\nu_8$ &
Spiritual assertion sharpens into focused cognition; strongly amplified.
\\
6B & $F_B$ & 6C & A &
$\nu_6,\nu_1$ &
Mental projection intensifies emotional drive; amplification.
\\
6C & $F_C$ & 6D & P &
$\nu_6,\nu_7$ &
Emotional assertiveness $\to$ physical action; preserved.
\\
6D & $F_D$ & 6D & P &
$\nu_6,\nu_4$ &
Terminal manifestation.
\\
\midrule

%----- Rhythm (7A–7D)
7A & $F_A$ & 7B & P &
$\nu_7,\nu_8$ &
Spiritual rhythm $\to$ mental-cycle recognition; preserved.
\\
7B & $F_B$ & 7C & P &
$\nu_7,\nu_1$ &
Thought cycles $\to$ emotional rhythmicity; preserved.
\\
7C & $F_C$ & 7D & W &
$\nu_7,\nu_2$ &
Emotional rhythms $\to$ physical biorhythms; weakened slightly.
\\
7D & $F_D$ & 7D & P &
$\nu_7,\nu_4$ &
Terminal physical rhythm system.
\\

\end{longtable}
\normalsize

\section{RPM Monad Interaction with Sub-Foundations}

\renewcommand{\arraystretch}{1.28}
\small
\begin{longtable}{c|c|c|p{2cm}|p{1.5cm}|p{4.5cm}}
\caption{RPM Monad Interaction with Sub-Foundations}
\label{tab:rpm-subfoundation} \\
\toprule
\textbf{Sub-F}
& \textbf{RPM Stage}
& \textbf{Outcome}
& \textbf{Composition}
& \textbf{Diag.\ Class}
& \textbf{Semantic Interpretation}
\\
\midrule
\endfirsthead
\multicolumn{6}{c}{\textit{Table \ref{tab:rpm-subfoundation} continued from previous page}} \\
\toprule
\textbf{Sub-F}
& \textbf{Stage}
& \textbf{Out}
& \textbf{Composition}
& \textbf{Diag.}
& \textbf{Interpretation}
\\
\midrule
\endhead
\midrule
\multicolumn{6}{r}{\textit{Continued on next page}} \\
\endfoot
\bottomrule
\endlastfoot

%---- 1 Awareness (A–D)
1A & $A0$ & Success &
$\text{unit} \gg= F_{1A}$ &
Clarity Init &
Spiritual awareness activates pure desire; RPM chain initializes cleanly.
\\
1B & $W_1$ & Delay &
$\text{RPM} \gg= F_{1B}$ &
Cog.\ Gate &
Mental clarity reveals hidden contradictions; pauses until resolved.
\\
1C & $P_1$ & Block &
$F_{1C} \gg= \text{Fail}$ &
Emot.\ Sat. &
Emotional awareness surfaces unresolved content; RPM halts.
\\
1D & $D_1$ & Success &
$F_{1D} \circ \text{unit}$ &
Embod.\ Anchor &
Physical awareness aligns intent with physical feasibility.
\\
\midrule

%---- 2 Positive (A–D)
2A & $A0$ & Success &
$\nu_2 \circ \text{unit}$ &
Desire Amp &
Positive spiritual stance increases motivational coherence.
\\
2B & $W_2$ & Success &
$\text{RPM} \gg= \nu_2$ &
Cog.\ Reinf &
Mental positivity supports stable belief structures in RPM.
\\
2C & $P_2$ & Delay &
$\nu_2 \gg= \text{hold}$ &
Emot.\ Reinf &
Emotional uplift softens blockages but may delay grounding.
\\
2D & $D_2$ & Success &
$\nu_2 \circ \text{ground}$ &
Somatic Int &
Physical health aligns with material acquisition.
\\
\midrule

%---- 3 Negative (A–D)
3A & $A0$ & Escalation &
$\nu_3 \circ \text{unit}$ &
Shadow Trig &
Spiritual negativity surfaces unresolved karmic material; RPM escalates.
\\
3B & $W_3$ & Block &
$F_{3B} \gg= \text{Fail}$ &
Cog.\ Distort &
Negative beliefs block RPM until reframed.
\\
3C & $P_3$ & Escalation &
$\nu_3 \gg= \nu_7$ &
Somatic Alarm &
Emotional turbulence intensifies physical resistance signals.
\\
3D & $D_3$ & Block &
$\text{halt}(\nu_3)$ &
Physio.\ Stress &
RPM halts due to somatic overload.
\\
\midrule

%---- 4 Vibration (A–D)
4A & $A0$ & Success &
$\nu_4 \circ \text{unit}$ &
Freq.\ Align &
Spiritual vibration enhances coherence of desire signal.
\\
4B & $W_4$ & Success &
$\text{RPM} \gg= \nu_4$ &
Cog.\ Reson &
Mental frequency stabilizes thought cycles.
\\
4C & $P_4$ & Success &
$\nu_4 \gg= \text{sync}$ &
Emot.\ Reson &
Emotional movement aligns with RPM pathway.
\\
4D & $D_4$ & Success &
$\text{ground} \circ \nu_4$ &
Energ.\ Embod &
Physical oscillation aligns with actionable motion.
\\
\midrule

%---- 5 Feminine (A–D)
5A & $A0$ & Delay &
$\nu_5 \gg= \text{unit}$ &
Recept.\ Gate &
Spiritual receptivity gathers missing prerequisites.
\\
5B & $W_5$ & Success &
$\text{RPM} \gg= \nu_5$ &
Intuit.\ Boost &
Mental openness enhances RPM's inference chain.
\\
5C & $P_5$ & Delay &
$\nu_5 \gg= \text{soften}$ &
Emot.\ Soften &
Emotional compassion reduces block intensity.
\\
5D & $D_5$ & Success &
$\nu_5 \circ \text{ground}$ &
Mat.\ Embrace &
Physical receptivity allows environmental alignment.
\\
\midrule

%---- 6 Masculine (A–D)
6A & $A0$ & Escalation &
$\nu_6 \circ \text{unit}$ &
Willforce Trig &
Spiritual will intensifies intent into force.
\\
6B & $W_6$ & Success &
$\text{RPM} \gg= \nu_6$ &
Cog.\ Drive &
Mental determination strengthens prerequisite chain.
\\
6C & $P_6$ & Success &
$\nu_6 \gg= \text{act}$ &
Emot.\ Impetus &
Will converts emotional pressure into direction.
\\
6D & $D_6$ & Success &
$\nu_6 \circ \text{ground}$ &
Phys.\ Action &
Physical execution ignites manifestation.
\\
\midrule

%---- 7 Rhythm (A–D)
7A & $A0$ & Delay &
$\nu_7 \circ \text{unit}$ &
Cycle Init &
Spiritual rhythmicity establishes iteration patterns in RPM.
\\
7B & $W_7$ & Success &
$\text{RPM} \gg= \nu_7$ &
Cycle Recog &
Cognitive cycles become explicit.
\\
7C & $P_7$ & Delay &
$\nu_7 \gg= \text{phase}$ &
Emot.\ Rhythm &
Emotional cycles phase-align with RPM progression.
\\
7D & $D_7$ & Success &
$\nu_7 \circ \text{ground}$ &
Phys.\ Sync &
Biorhythms align with material timing.
\\

\end{longtable}
\normalsize

%% ============================================================================
%% CHAPTER 8: ADVANCED CLASSIFICATION AND FORMAL SYSTEM TABLES
%% ============================================================================

\chapter{Advanced Classification and Formal System Tables}
\label{ch:advanced-tables}

\begin{tcolorbox}[colback=blue!5!white,colframe=blue!75!black,title=\vnewfour,breakable]
This chapter provides advanced classification tables for fractal stability,
transfinite convergence, multi-goal RPM interactions, compiler mappings,
and category-theoretic structure summaries.
\end{tcolorbox}

\section{Fractal Stability Classes vs Sub-Foundations}

\renewcommand{\arraystretch}{1.3}
\small
\begin{longtable}{c|c|c|p{2cm}|p{5cm}}
\caption{Fractal Stability Classes vs Sub-Foundations}
\label{tab:fractal-stability-subfoundation} \\
\toprule
\textbf{Sub-F}
& \textbf{Stability}
& \textbf{Behavior}
& \textbf{Convergence}
& \textbf{Interpretation}
\\
\midrule
\endfirsthead
\multicolumn{5}{c}{\textit{Table \ref{tab:fractal-stability-subfoundation} continued from previous page}} \\
\toprule
\textbf{Sub-F}
& \textbf{Stability}
& \textbf{Behavior}
& \textbf{Convergence}
& \textbf{Interpretation}
\\
\midrule
\endhead
\midrule
\multicolumn{5}{r}{\textit{Continued on next page}} \\
\endfoot
\bottomrule
\endlastfoot

1A & Stable & Fixed & Strong &
Spiritual awareness forces fractal to a stable fixed point. \\

1B & Semi-Stable & Oscillatory & Weak &
Cognitive reflection induces bounded oscillation then stabilization. \\

1C & Unstable & Divergent & None &
Emotional saturation destabilizes fractal until emotional block is resolved. \\

1D & Stable & Fixed & Strong &
Physical grounding reduces fractal drift. \\
\midrule

2A & Stable & Contractive & Strong &
Positive spiritual influence tightens fractal mapping. \\

2B & Stable & Fixed & Strong &
Mental positivity produces calm attractor states. \\

2C & Semi-Stable & Exp$\to$Contr & Conditional &
Emotionally positive states amplify then settle. \\

2D & Stable & Fixed & Immediate &
Physical health signals produce stable patterns. \\
\midrule

3A & Unstable & Divergent & None &
Shadow activation destabilizes fractal mapping. \\

3B & Unstable & Chaotic & None &
Negative cognition introduces chaotic jitter. \\

3C & Unstable & Chaotic$\to$Exp & None &
Emotional turmoil keeps fractal from stabilizing. \\

3D & Unstable & Expansive & None &
Physiological stress prevents any attractor formation. \\
\midrule

4A & Stable & Periodic & Conv.\ Cycle &
Spiritual vibration introduces rhythmic stabilization. \\

4B & Stable & Periodic & Conv.\ Cycle &
Mental rhythm reveals convergence through periodicity. \\

4C & Semi-Stable & Period$\to$Fixed & Conditional &
Emotional vibrational field can stabilize oscillations. \\

4D & Stable & Fixed & Strong &
Physical vibration locks into frequency alignment. \\

\end{longtable}
\normalsize

\section{Transfinite Fractal Convergence Classes}

\renewcommand{\arraystretch}{1.25}
\small
\begin{longtable}{c|p{2cm}|p{2cm}|p{2.5cm}|p{4.5cm}}
\caption{Transfinite Fractal Convergence Classes}
\label{tab:transfinite-convergence} \\
\toprule
\textbf{Class}
& \textbf{Ordinal Domain}
& \textbf{Limit Behavior}
& \textbf{Evaluation Form}
& \textbf{Interpretation}
\\
\midrule
\endfirsthead
\multicolumn{5}{c}{\textit{Table \ref{tab:transfinite-convergence} continued from previous page}} \\
\toprule
\textbf{Class}
& \textbf{Ordinal}
& \textbf{Limit}
& \textbf{Evaluation}
& \textbf{Interpretation}
\\
\midrule
\endhead
\midrule
\multicolumn{5}{r}{\textit{Continued on next page}} \\
\endfoot
\bottomrule
\endlastfoot

$\mathsf{C}_0$ & Finite ($n<\omega$) & Fixed & $\sem{\Phi_n}(x)$ &
Finite fractals always terminate in fixed points. \\

$\mathsf{C}_1$ & $\omega$ & Contractive Limit & $\bigsqcup_{n<\omega}\sem{\Phi_n}(x)$ &
Infinite iteration yields least fixed point. \\

$\mathsf{C}_2$ & $\omega \cdot k$ & Periodic$\to$Fixed & Limit over $k$ cycles &
Ordinal blocks collapse to stable attractor. \\

$\mathsf{C}_3$ & $\omega^2$ & Layered Limit & Double-structured limit &
Fractal converges in hierarchical layers. \\

$\mathsf{C}_4$ & $\omega^\omega$ & Self-similar Limit & Transfinite recursion &
Deep self-similarity yields canonical attractor form. \\

\end{longtable}
\normalsize

\section{RPM Multi-Goal Interaction Matrix}

\renewcommand{\arraystretch}{1.25}
\small
\begin{longtable}{c|c|c|c|p{5cm}}
\caption{RPM Multi-Goal Interaction Matrix}
\label{tab:rpm-multigoal} \\
\toprule
\textbf{Goal Pair}
& \textbf{Dependency}
& \textbf{RPM Outcome}
& \textbf{Conflict}
& \textbf{Interpretation}
\\
\midrule
\endfirsthead
\multicolumn{5}{c}{\textit{Table \ref{tab:rpm-multigoal} continued from previous page}} \\
\toprule
\textbf{Pair}
& \textbf{Dep.}
& \textbf{Outcome}
& \textbf{Conflict}
& \textbf{Interpretation}
\\
\midrule
\endhead
\midrule
\multicolumn{5}{r}{\textit{Continued on next page}} \\
\endfoot
\bottomrule
\endlastfoot

(F1, F2) & Cooperative & Merge & None &
Unity and Wisdom reinforce each other. \\

(F2, F6) & Semi-Coop & Merge$\to$Split & Low &
Wisdom supports Material, but may force reprioritization. \\

(F3, F7) & Oppositional & Split & High &
Life and Lust compete for energy allocation. \\

(F4, F5) & Hybrid & Merge$\to$Pause$\to$Merge & Medium &
Companionship alternates between support and emotional delay. \\

(F6, F1) & Cooperative & Merge & None &
Material flows support Unity expression. \\

\end{longtable}
\normalsize

\section{ACBE Functor Interaction Table}

\renewcommand{\arraystretch}{1.25}
\small
\begin{longtable}{c|c|p{2.2cm}|p{2.5cm}|p{4.5cm}}
\caption{ACBE Functor Interaction Table}
\label{tab:acbe-interaction} \\
\toprule
\textbf{World Level}
& \textbf{Functor}
& \textbf{Preservation}
& \textbf{Mapping Form}
& \textbf{Interpretation}
\\
\midrule
\endfirsthead
\multicolumn{5}{c}{\textit{Table \ref{tab:acbe-interaction} continued from previous page}} \\
\toprule
\textbf{Level}
& \textbf{Functor}
& \textbf{Preservation}
& \textbf{Mapping}
& \textbf{Interpretation}
\\
\midrule
\endhead
\midrule
\multicolumn{5}{r}{\textit{Continued on next page}} \\
\endfoot
\bottomrule
\endlastfoot

A$\to$B & $F_B$ & Structure-preserving & $A_n^{k} \mapsto B_n^{k}$ &
Spiritual $\to$ Mental translation preserves operator semantics. \\

B$\to$C & $F_C$ & Rhythm-preserving & $B_n^{k} \mapsto C_n^{k}$ &
Cognition $\to$ Emotion introduces dynamic oscillation. \\

C$\to$D & $F_D$ & Ground-preserving & $C_n^{k} \mapsto D_n^{k}$ &
Emotion $\to$ Physical anchors noetic transformation. \\

A$\to$D & $F_D \circ F_C \circ F_B$ & Fully functorial & $A_n^{k} \mapsto D_n^{k}$ &
Full ACBE descent preserves global semantics. \\

\end{longtable}
\normalsize

\section{Noetic Operator Spectrum Table}

\renewcommand{\arraystretch}{1.25}
\small
\begin{longtable}{c|c|c|p{2cm}|p{5cm}}
\caption{Noetic Operator Spectrum Table}
\label{tab:noetic-spectrum} \\
\toprule
\textbf{Noetic}
& \textbf{Spectral Radius}
& \textbf{Eigenclass}
& \textbf{Fixed-Point}
& \textbf{Interpretation}
\\
\midrule
\endfirsthead
\multicolumn{5}{c}{\textit{Table \ref{tab:noetic-spectrum} continued from previous page}} \\
\toprule
\textbf{Noetic}
& \textbf{Radius}
& \textbf{Eigenclass}
& \textbf{Fixed-Pt}
& \textbf{Interpretation}
\\
\midrule
\endhead
\midrule
\multicolumn{5}{r}{\textit{Continued on next page}} \\
\endfoot
\bottomrule
\endlastfoot

$\nu_0$ & 0 & Identity & All fixed & Neutral operator. \\

$\nu_1$ & 1 & Awareness & Stable fixed & Mind stabilizes itself. \\

$\nu_2$ & $<1$ & Contractive & Strong fixed & Positive shrinks deviation. \\

$\nu_3$ & $>1$ & Expansive & None & Negative amplifies error. \\

$\nu_4$ & 1 & Oscillatory & Periodic & Vibration induces cycles. \\

$\nu_5$ & $<1$ & Contractive & Stable & Feminine integrates. \\

$\nu_6$ & $>1$ & Expansive & Unstable & Masculine forces expansion. \\

$\nu_7$ & 1 & Rhythmic & Periodic & Repetition introduces harmonic motion. \\

$\nu_8$ & $<1$ & Convergent & High-fixed & Above elevates. \\

$\nu_9$ & 1 & Anchoring & Ground-fixed & Below grounds. \\

\end{longtable}
\normalsize

\section{Eigenfractal Signature Table}

\renewcommand{\arraystretch}{1.25}
\small
\begin{longtable}{p{2.5cm}|c|p{2cm}|c|p{5cm}}
\caption{Eigenfractal Signature Table}
\label{tab:eigenfractal-signature} \\
\toprule
\textbf{Fractal}
& \textbf{Eigenvalue}
& \textbf{Fixed-Point Class}
& \textbf{Stability}
& \textbf{Interpretation}
\\
\midrule
\endfirsthead
\multicolumn{5}{c}{\textit{Table \ref{tab:eigenfractal-signature} continued from previous page}} \\
\toprule
\textbf{Fractal}
& \textbf{Eigen}
& \textbf{FP Class}
& \textbf{Stab.}
& \textbf{Interpretation}
\\
\midrule
\endhead
\midrule
\multicolumn{5}{r}{\textit{Continued on next page}} \\
\endfoot
\bottomrule
\endlastfoot

$\langle 1:4:7\rangle$ & 1 & Periodic & Stable &
MVR produces rhythmic awareness cycles. \\

$\langle 2:3\rangle$ & 0 & Neutralized & Stable &
Dual polarity collapses to identity. \\

$\langle 5:6:1\rangle$ & 1 & Fixed & Semi-Stable &
Gender integration loops into awareness. \\

$\langle 8:1:4:7:9\rangle$ & 1 & Fixed & Strong &
ACBE descent is a stable global attractor. \\

\end{longtable}
\normalsize

\section{Monad Law Verification Table}

\renewcommand{\arraystretch}{1.25}
\small
\begin{longtable}{c|p{4cm}|p{6.5cm}}
\caption{RPM Monad Law Verification}
\label{tab:rpm-monad-laws} \\
\toprule
\textbf{Law}
& \textbf{Form}
& \textbf{Interpretation}
\\
\midrule
\endfirsthead
\multicolumn{3}{c}{\textit{Table \ref{tab:rpm-monad-laws} continued from previous page}} \\
\toprule
\textbf{Law}
& \textbf{Form}
& \textbf{Interpretation}
\\
\midrule
\endhead
\midrule
\multicolumn{3}{r}{\textit{Continued on next page}} \\
\endfoot
\bottomrule
\endlastfoot

Left Identity & $\text{unit}(x) \gg= f = f(x)$ &
Pure desire correctly seeds the monadic chain. \\

Right Identity & $m \gg= \text{unit} = m$ &
RPM preserves the evolved result after completion. \\

Associativity & $(m \gg= f) \gg= g = m \gg= (\lambda x.\, f(x) \gg= g)$ &
Sequential prerequisite checking is structurally coherent. \\

\end{longtable}
\normalsize

\section{TKS Runtime Error and Exception Classes}

\renewcommand{\arraystretch}{1.25}
\small
\begin{longtable}{c|p{2.5cm}|p{7cm}}
\caption{TKS Runtime Error \& Exception Classes}
\label{tab:tks-errors} \\
\toprule
\textbf{Code}
& \textbf{Name}
& \textbf{Interpretation}
\\
\midrule
\endfirsthead
\multicolumn{3}{c}{\textit{Table \ref{tab:tks-errors} continued from previous page}} \\
\toprule
\textbf{Code}
& \textbf{Name}
& \textbf{Interpretation}
\\
\midrule
\endhead
\midrule
\multicolumn{3}{r}{\textit{Continued on next page}} \\
\endfoot
\bottomrule
\endlastfoot

E01 & Missing Desire & A0 not satisfied; RPM cannot begin. \\

E02 & Foundation Block & A D/W/P prerequisite fails. \\

E03 & Fractal Divergence & Fractal evaluation does not converge. \\

E04 & Type Mismatch & Ill-typed TKS expression. \\

E05 & Category Violation & Functor mapping inconsistent across worlds. \\

\end{longtable}
\normalsize

\section{Type System Subtyping Lattice}

\renewcommand{\arraystretch}{1.25}
\small
\begin{longtable}{p{2cm}|p{2.5cm}|p{6.5cm}}
\caption{Type System Subtyping Lattice}
\label{tab:type-lattice} \\
\toprule
\textbf{Subtype}
& \textbf{Supertype}
& \textbf{Meaning}
\\
\midrule
\endfirsthead
\multicolumn{3}{c}{\textit{Table \ref{tab:type-lattice} continued from previous page}} \\
\toprule
\textbf{Subtype}
& \textbf{Supertype}
& \textbf{Meaning}
\\
\midrule
\endhead
\midrule
\multicolumn{3}{r}{\textit{Continued on next page}} \\
\endfoot
\bottomrule
\endlastfoot

Aspect & Foundation & Every aspect belongs to one foundation. \\

Foundation & Domain & All foundations inhabit Idea-space. \\

Noetic & Domain$\to$Domain & Operators act on Idea-space. \\

Fractal & Domain$\to$Domain & Higher-order operator. \\

RPM[$\tau$] & $\tau$ & Monadic wrapper preserves type. \\

\end{longtable}
\normalsize

\section{Operational Semantics Reduction Table}

\renewcommand{\arraystretch}{1.25}
\small
\begin{longtable}{p{3.5cm}|p{8cm}}
\caption{Operational Semantics Reduction Rules}
\label{tab:operational-semantics} \\
\toprule
\textbf{Reduction}
& \textbf{Meaning}
\\
\midrule
\endfirsthead
\multicolumn{2}{c}{\textit{Table \ref{tab:operational-semantics} continued from previous page}} \\
\toprule
\textbf{Reduction}
& \textbf{Meaning}
\\
\midrule
\endhead
\midrule
\multicolumn{2}{r}{\textit{Continued on next page}} \\
\endfoot
\bottomrule
\endlastfoot

$E^{k} \to \nu_k(E)$ &
Noetic application reduces exponent to operator application. \\

$E_1 \oplus_T E_2 \to E_1 \cup E_2$ &
Tootra-addition reduces to union. \\

$E_1 \ominus_T E_2 \to E_1 \setminus E_2$ &
Tootra-subtraction reduces to set difference. \\

$\langle k_1:k_2 \rangle(E) \to \nu_{k_1}(\nu_{k_2}(E))$ &
Fractal reduction unfolds chaining. \\

$\text{RPM}(x) \to \text{Failure}(A_m)$ &
RPM aborts when prerequisite fails. \\

\end{longtable}
\normalsize

\section{Denotational Semantics Domain Table}

\renewcommand{\arraystretch}{1.25}
\small
\begin{longtable}{c|p{2.5cm}|p{7cm}}
\caption{Denotational Semantics Domains}
\label{tab:denotational-domains} \\
\toprule
\textbf{Symbol}
& \textbf{Domain}
& \textbf{Meaning}
\\
\midrule
\endfirsthead
\multicolumn{3}{c}{\textit{Table \ref{tab:denotational-domains} continued from previous page}} \\
\toprule
\textbf{Symbol}
& \textbf{Domain}
& \textbf{Meaning}
\\
\midrule
\endhead
\midrule
\multicolumn{3}{r}{\textit{Continued on next page}} \\
\endfoot
\bottomrule
\endlastfoot

$\sem{E}$ & $\mathcal{D}$ & Denotation of expression. \\

$\mathcal{D}_\bot$ & Domain + bottom & Partial execution domain. \\

$\text{Noe}$ & $\mathcal{D}\to\mathcal{D}$ & Noetic operators. \\

$\text{Frac}$ & $(\mathcal{D}\to\mathcal{D})$ & Fractals as higher operators. \\

$\text{RPM}(\tau)$ & $\tau + \mathsf{Error}$ & Effect monad. \\

\end{longtable}
\normalsize

\section{Canonical Equivalence Table}

\renewcommand{\arraystretch}{1.25}
\small
\begin{longtable}{p{2.5cm}|p{3.5cm}|p{5.5cm}}
\caption{Canonical Equivalence Relations in TKS}
\label{tab:canonical-equivalence} \\
\toprule
\textbf{Form}
& \textbf{Equivalent Form}
& \textbf{Meaning}
\\
\midrule
\endfirsthead
\multicolumn{3}{c}{\textit{Table \ref{tab:canonical-equivalence} continued from previous page}} \\
\toprule
\textbf{Form}
& \textbf{Equivalent}
& \textbf{Meaning}
\\
\midrule
\endhead
\midrule
\multicolumn{3}{r}{\textit{Continued on next page}} \\
\endfoot
\bottomrule
\endlastfoot

$E^{23}$ & $E^{0}$ &
Dual collapse: $\nu_2\circ\nu_3 = \nu_0$. \\

$\langle 2:3 \rangle(E)$ & $E$ &
Polarity fractal neutralizes. \\

$A_n^{k} \mapsto D_n^{k}$ & $F_D(F_C(F_B(A_n^{k})))$ &
ACBE descent equivalence. \\

$\text{RPM}(x) = x$ & $\forall m.\, \text{Satisfied}(m)$ &
RPM identity if all prerequisites met. \\

\end{longtable}
\normalsize

\section{Compiler Mapping Table}

\renewcommand{\arraystretch}{1.25}
\small
\begin{longtable}{p{1.8cm}|p{2.2cm}|p{2.5cm}|p{5cm}}
\caption{Compiler Mapping Table: AST $\to$ IR $\to$ Runtime}
\label{tab:compiler-mapping} \\
\toprule
\textbf{AST Node}
& \textbf{IR Form}
& \textbf{Runtime Action}
& \textbf{Meaning}
\\
\midrule
\endfirsthead
\multicolumn{4}{c}{\textit{Table \ref{tab:compiler-mapping} continued from previous page}} \\
\toprule
\textbf{AST}
& \textbf{IR}
& \textbf{Runtime}
& \textbf{Meaning}
\\
\midrule
\endhead
\midrule
\multicolumn{4}{r}{\textit{Continued on next page}} \\
\endfoot
\bottomrule
\endlastfoot

Element & CONST(w,n) & Load value & Basic TKS element. \\

Power & APPLY(k) & Apply noetic & Exponent semantics. \\

Fractal & CHAIN(\ldots) & Unfold chain & Higher-order operator. \\

RPM & MONAD(seq) & Evaluate chain & Prerequisite engine. \\

Binary Op & OP(op) & Apply set-op & Tootra algebra. \\

\end{longtable}
\normalsize

\section{Category-Theoretic Structure Summary}

\renewcommand{\arraystretch}{1.25}
\small
\begin{longtable}{p{2cm}|p{2.5cm}|p{7cm}}
\caption{Category-Theoretic Structure Summary}
\label{tab:category-summary} \\
\toprule
\textbf{Structure}
& \textbf{Type}
& \textbf{Meaning}
\\
\midrule
\endfirsthead
\multicolumn{3}{c}{\textit{Table \ref{tab:category-summary} continued from previous page}} \\
\toprule
\textbf{Structure}
& \textbf{Type}
& \textbf{Meaning}
\\
\midrule
\endhead
\midrule
\multicolumn{3}{r}{\textit{Continued on next page}} \\
\endfoot
\bottomrule
\endlastfoot

Noetica & Category & Objects: Ideas; Morphisms: Noetics. \\

ACBE & Functor & World-lowering mapping. \\

RPM & Monad & Sequential prerequisite evaluation. \\

Fractals & Endofunctors & Higher-order noetic transformations. \\

Worlds & 2-Category & World transitions as 2-morphisms. \\

\end{longtable}
\normalsize
